\appendix

\noindent\textbf{Online Appendix.} The appendices that follow are for online publication. The run ledger of superseded figures and the pre-committed-rule adjudication record form a standalone document in the replication package (\texttt{replication\_appendices}). Its cross-references resolve against this manuscript.

\section{Agent-Based Model: Specification}\label{sec:method-abm}

I adopt an agent-based approach because the phenomenon under study is heterogeneous. Households differ in coupon, loan age, and exposure to search and sentiment frictions. Those differences interact non-linearly with the level of interest rates. \citet{geanakoplos2012} argue for agent-based models of mortgage prepayment in environments where behavioral and search frictions vary across the population and over time, and demonstrate their feasibility. That property is directly relevant here: the frictions I model (housing search costs and consumer sentiment) are themselves time-varying. \citet{tesfatsion2006} sets out the broader methodological case for agent-based computational economics in settings with meaningful heterogeneity.

The ABM simulates a population of N = 10,000 heterogeneous households. Initial agent endowments (income, home value) are drawn from log-normal distributions calibrated to Federal Reserve Economic Data (FRED) medians. In the production multi-vintage specification of Appendix~\ref{sec:method-frictions}, each household's mortgage coupon and origination date are drawn from the Federal Reserve's actual holdings composition. That composition is the live CUSIP-level SOMA parse of Section~\ref{sec:abm-results}, across seven coupon cohorts. A single-cohort variant serves exposition and mechanism analysis. It instead endows every agent with a 30-year fixed-rate mortgage at a 3.0\% coupon aged 60 months at the simulation start. Where the text below discusses a single representative coupon, it refers to this pedagogical case, not to the production population.

Section~\ref{sec:robustness} reports a scale-sensitivity replication showing that the gap in explanatory power is not an artifact of sample size. Whether it reflects modeling paradigm or data source is taken up in Section~\ref{sec:robustness-crossdesign}.

The decision to relocate weighs the utility of moving against the aggregate financial penalty of doing so. The monthly payment follows the standard annuity formula
\begin{equation}\label{eq:annuity}
M = P\,\frac{r\,(1+r)^{n}}{(1+r)^{n}-1},
\end{equation}
where $P$ is the outstanding principal, $r$ is the monthly interest rate, and $n$ is the remaining term in months. The mobility penalty evaluates \eqref{eq:annuity} twice. The first evaluation applies the borrower's locked-in note rate to the original principal and term, which gives the current payment. By the amortization identity, that payment equals the payment on the outstanding balance at the note rate over the remaining term. The second applies the prevailing market rate to the outstanding balance over a fresh full original term, which gives the payment the borrower would take on by refinancing. The replacement loan re-amortizes---a term-preference persistence the code implements deliberately. The penalty scales with the difference between the two. A native short-term variant of the penalty, active only for terms at or below 240 months, floors payment decreases at zero; Appendix~\ref{app:abmspec} specifies it with the full decision rule. Here $r$ is expressed in decimal monthly terms: the annual note rate divided by twelve, so that a 3.0\% coupon gives $r = 0.0025$. It is a rate \emph{level}, distinct in both role and units from the annualized rate \emph{gaps} $g_{i,t}$ (in 100 bp units) and $\tilde g_{i,t} = g_{i,t}/100$ (decimal) that drive the hazard framework of Section~\ref{sec:pathb}. No rate symbol is reused across the two. I evaluate the population's move/stay decisions household-by-household at each point of a three-dimensional grid of market rate, dynamic friction, and rate velocity (and, in the multi-vintage specification, across seven coupon cohorts). I then aggregate them to a pool-level annualized Conditional Prepayment Rate (CPR) surface. CPR is a pool-level convention. The household-level object is an annualized move intensity, consumed by the accounting layer at one-twelfth per month (Appendix~\ref{app:abmspec}). That surface is precomputed once, then interpolated month-by-month against the live population of simulated market conditions. That cuts the cost of the full grid by roughly two orders of magnitude and makes fifty-seed Monte Carlo estimation feasible. The surface stores each node's direct household-by-household evaluation, so interpolation reproduces the micro-simulation exactly at grid points. The operative approximation is the trilinear segment between 50 bp-spaced nodes. A spot-check validates that approximation directly. After an on-grid parity gate reproduced the committed surface values, direct micro-simulation and interpolation at worst-case off-grid midpoints differ by at most 0.23 points of CPR away from the freeze threshold, with the Danish column exact. The one systematic difference is the wait-and-see freeze itself. The engine applies it as a step at the 150 bp trigger, while the surface renders it as a ramp across the adjacent velocity nodes. That gap binds in four realized months (August--November 2022; SOMA-weighted mean absolute difference 0.22 points of CPR over the window). These are surface semantics that the production figures share. The binary-step leg of Appendix~\ref{sec:robustness-extensions}'s freeze sweep prices their trapped-liquidity consequence directly. Figure~\ref{fig:cprsurface} renders one slice of that surface (the 3.0\% cohort of the expository variant) with the realized QT window located on it.

\begin{figure}[H]
\centering
\includegraphics[width=0.85\textwidth]{fig8_cpr_surface}
\caption{The precomputed CPR surface: the production surface's slice at the 3.0\% thirty-year SOMA cohort, at its parsed seasoning. Left: the U.S. par-payoff surface over the market-rate $\times$ dynamic-friction grid (no-freeze regime), with the 42 QT-window months overlaid at their realized (rate, friction) coordinates; the window sits entirely in the deep out-of-the-money corner, near the baseline turnover floor. Right: rate slices at the friction grid point nearest the window mean. Velocity enters the grid only through the wait-and-see freeze of Appendix~\ref{sec:method-frictions} (20\% of otherwise-mobile households held in place above a 150 bp trailing six-month rise), so the two regimes shown are its entire velocity structure; the Berger-recalibrated Danish counterfactual (Section~\ref{sec:abm-danish}) is near-flat, moving roughly a quarter point across the grid. The U.S. slice is identical across the frozen production surface and the current code head; the Danish curve is the Berger production specification.}
\label{fig:cprsurface}
\end{figure}

\begin{figure}[H]
\centering
\includegraphics[width=0.9\textwidth]{fig10_abm_cpr_paths}
\caption{Production ABM simulated CPR against the empirical SOMA back-out over the QT window (frozen fold-in run, seed 42). The simulated path runs at roughly twice the empirical level and does not track its month-to-month shape; the annotated means, lag-0 $r$, and $R^2$ are the manifest values quoted in the text. This is the level-and-shape mismatch behind the negative path statistics; the cross-correlation structure across lags is in Figure~\ref{fig:ccf}.}
\label{fig:abmpaths}
\end{figure}

\begin{figure}[H]
\centering
\includegraphics[width=0.82\textwidth]{monte_carlo_trapped_liquidity}
\caption{Monte Carlo population-draw distribution of ABM trapped liquidity (50 seeds, fold-in specification; CPR surface rebuilt per draw). The seed-level dispersion (SD \$24.5 billion) dwarfs the distance between the frozen seed-42 production draw (\$91.0 billion, 34th percentile) and the seed mean (\$103.7 billion); the empirical benchmark of \$764.7 billion lies far outside the distribution's support, which is the visual form of the ABM's 10--17\% recovery band.}
\label{fig:mc}
\end{figure}

The simulation contrasts two institutional payoff rules. Under U.S. par-payoff, the cost to retire the mortgage is strictly equal to the outstanding principal. This is the negatively convex case: the investor's upside is capped, while extension risk is bounded only by the scheduled-amortization profile. That profile runs to a 14.7-year weighted-average life at zero CPR against the no-shock 3.4-year life (Table~\ref{tab:wal}), an order-of-magnitude extension ceiling. Under Danish market-value repurchase, the payoff cost is the present value of the remaining scheduled payments discounted at the current 30-year primary mortgage rate (FRED MORTGAGE30US), capped at par. Two qualifications attach to this proxy. First, it is an option-free present value discounted at the primary rate, so it ignores the embedded par call and any option-adjusted spread. The primary rate also embeds the primary--secondary spread (guarantee fee, servicing, originator margin; \citealp{fuster2013}), whereas the bond-market discounting a true buyback would use sits below it. Both choices overstate the borrower's buyback discount as rates rise. These biases inflate the mechanism-substitution Danish CPR of Section~\ref{sec:abm-danish} (the since-superseded 47.1\%), but they are exactly immaterial to the Berger-recalibrated production finding. Under neither transplant anchor does that leg's prepayment hazard consume this pricing proxy: the Danish-level anchor imports the near-flat moving curve, and the production U.S.-intercept anchor evaluates the U.S. hazard at zero rate gap. A replication-package diagnostic re-runs the Danish leg with the discount rate shifted down by 25--150 bp. It moves trapped liquidity by \$0.00 billion. It flips 2,470--92,495 of the 1,683,082 loan-month present-value classifications\footnote{This Danish-leg classification universe of 1{,}683{,}082 loan-months is 42 fewer than the 1{,}683{,}124 evaluated in the U.S.\ production hazard panel (Section~\ref{sec:pathb}): the two figures count active loan-months in the two counterfactual legs, whose prepayment paths differ marginally even under the shared imported elasticity,} that no downstream quantity in that leg consumes. Second, what the Danish rule neutralizes is extension risk for the bondholder \citep{berger2026}. The borrower retains the par call, so the rate-decline side keeps its price cap, and I use ``convexity-neutral'' below strictly in this rate-rise sense. Under the Danish rule, as market rates rise above 3.0\%, the bond's present value falls. That lets the borrower extinguish the debt at a discount and guarantees continuous cash flow generation. I simulate the Danish counterfactual balance forward on a dynamic, declining-balance basis: each month's elevated Danish roll-off applies to the already-shrunk balance carried over from the prior month, so the simulated payoff path compounds correctly over the multi-year QT window.

\section{Macroeconomic Friction and Structural Extensions}\label{sec:method-frictions}

To quantify the behavioral friction gap in standard models, I scale base transaction friction using two FRED series. (1) I measure search friction as the shortfall in Active Listing Count (ACTLISCOUUS) against a 2017--2019 baseline. That window is the last stable pre-pandemic period, before the 2020/21 rate-driven frenzy and the 2022 inventory collapse. (2) I index psychological friction from University of Michigan Consumer Sentiment (UMCSENT) against a baseline of 85.0. Readings below that line score as incremental friction at +5 bp per point short, capped at +150 bp during acute sentiment stress. The result is a dynamic friction parameter---a one-time moving transaction cost expressed as a share of the household's current home value. It scales between 8.23\% and 9.96\% of home value over the 42 active QT months.

I add three behavioral extensions to the household decision function. No parameter in any of them is tuned to match the empirical benchmark. Their grounding is nonetheless uneven: the first two are parameterized from external sources, the third is an author calibration. The first is a debt-to-income (DTI) constraint that vetoes any move in which the new housing payment would exceed 43\% of gross monthly income. That threshold is a front-end heuristic motivated by the pre-2021 General QM debt-to-income limit (12 C.F.R. \S{} 1026.43(e)(2)(vi)). The regulatory threshold was a back-end total-debt ratio, and the price-based General QM definition superseded it in 2021, so I use it as a behavioral anchor. The model carries no total-debt data against which to test the threshold (Section~\ref{sec:limitations}). The second is asymmetric loss aversion that scales perceived payment increases by 2.25$\times$ while leaving decreases unscaled (loss aversion: \citealp{kahneman1979}; the 2.25 coefficient: \citealp{tversky1992}). \citepos{andersen2022} housing-market estimate of roughly 2.5 is close. Any coefficient in this range triggers the same turnover-floor recalibration described below, so the artifact interpretation of Section~\ref{sec:abm-results} does not hinge on the choice. The third is a wait-and-see freeze that holds 20\% of otherwise-mobile households in place whenever the trailing six-month change in mortgage rates exceeds 150 bp. It is modeled as a fixed per-agent trait, so the CPR surface remains deterministic. The 20\% share and 150 bp trigger are author calibrations with no external source. So are the sentiment-to-friction mapping and the base transaction-friction level of the preceding paragraph: the UMCSENT and ACTLISCOUUS inputs are data, but the $+5$ bp-per-point mapping, its $+150$ bp cap, and the base friction level are mine. Section~\ref{sec:robustness-crossdesign} tests an externally parameterized variant of the mobility-desire scale and the freeze share, drawn from the same source as Path B's elasticity.

Loss aversion steepens the penalty function facing every household. The model's mobility-desire parameter is independently calibrated to reproduce a real-world baseline turnover floor of 4--5\% CPR at 8\% market rates. That floor is the turnover speed observed on deeply out-of-the-money discount cohorts during 2023--2024. Refinancing incentives were absent then, and remaining prepayment reflected death, divorce, and forced relocation. I therefore hold that empirical floor fixed and re-run the calibration under the steeper penalty function. The re-run requires nearly tripling the mobility-desire scale parameter. That mechanical recalibration is the principal driver of the results in Section~\ref{sec:abm-results}.

Whenever a specification change alters the penalty or covariate structure facing every household, the mobility-desire scale is re-derived --- that parameter and no other --- so the simulation still reproduces the empirically observed 4--5\% baseline turnover floor. The same rule generates the recalibrated variant of the cross-design test (Section~\ref{sec:robustness-crossdesign}); suspending it defines the frozen variant. Appendix~\ref{sec:robustness-extensions} reports how much work this single choice does. Because the 4--5\% floor is imposed by construction, the ABM's agreement with observed discount-cohort turnover at deep rate gaps is an input, not a finding. The model earns no evidential credit on that dimension. Its evidential exposure lies entirely where the anchor leaves it free: the dollar recovery and the timing of the CPR path. The anchor is also insulated from the behavioral gates by construction. The calibration point is evaluated at zero rate velocity, where the wait-and-see freeze cannot bind, so no freeze parameterization can move the floor; Appendix~\ref{sec:robustness-extensions} verifies that the calibrated scale and floor CPR are bit-identical at every freeze share swept. What the anchor purchases, and what dropping it costs, is measured directly in Section~\ref{sec:robustness-crossdesign}. The externally parameterized variant that removes the floor overshoots the benchmark by half, while predicting essentially zero turnover at QT-window rate gaps.

I added an income-scaled curtailment channel as an extension specified ex ante. The ABM's monthly CPR over-prediction correlates with year-over-year growth in real disposable income (FRED series DSPIC96) at $r = -0.487$. The correlation concentrates in the worst over-prediction months, June--September 2022. That diagnostic led me to predict ex ante that curtailment, a strictly additive channel, would worsen the aggregate match. I implemented it as an income-scaled addition to simulated roll-off. The mapping is linear in real income growth: zero curtailment at $-2$\% growth, rising to a full-scale annualized rate of 1.2\% CPR at +4\% growth, consistent with GSE daily prepayment report estimates.

Finally, as a diagnostic, I replaced the single representative 3.0\% coupon assumed elsewhere with a multi-vintage composition. Bucketing SOMA holdings into seven coupon cohorts yields a weighted-average coupon of approximately 2.55\% on the 30-year book, with the 2.0\% and 2.5\% buckets accounting for roughly 69\% of 30-year face value. A lower coupon implies a larger gap to prevailing market rates, so my ex ante hypothesis was that this composition would increase predicted trapped liquidity toward the empirical benchmark. This pass excludes fifteen-year MBS (approximately 9\% of SOMA face value), and cohort origination dates are approximated from parsed maturity strings under a 360-month term assumption. Section~\ref{sec:abm-results} describes the production replacement of this diagnostic with a live, CUSIP-level pull of SOMA holdings.

\section{ABM Decision Rule and Parameter Provenance}\label{app:abmspec}
% Sources: abm/abm_lockin_simulation.py ("S:"), abm/fed_mbs_extension_risk.py ("F:"),
% abm/cross_design_test.py, abm/abm_external_gates.py. Line pins refer to the committed code.

This appendix specifies the ABM's move decision and every behavioral parameter in the detail needed to reimplement the model without consulting the code. Each equation and value is pinned to its source line in the replication package.

\paragraph{The move decision.}
Let $A(P,r,n) = P\,(r/12)\big/\big(1-(1+r/12)^{-n}\big)$ be the monthly payment on principal $P$ at annual rate $r$ over $n$ months. % S:130-136
Household $i$ carries gross income $y_i$ (\$/year), home value $H_i$ (\$), original principal $P_i = 0.8\,H_i$ (a fixed 80\% origination LTV in the synthetic population), note rate $c_i$, original term $n_i$ in months, and outstanding balance $K_i$ (closed form at the cohort's seasoning; the U.S. par payoff). % S:141-155, S:314
The replacement-loan payment is $M^{\mathrm{new}}_{i,t} = A(K_i, m_t, n_i)$: the par payoff financed at the market rate $m_t$ over a fresh \emph{full original term}, not the remaining term. That is term-preference persistence. % S:222-225, S:366-381
With $\Delta_{i,t} = M^{\mathrm{new}}_{i,t} - A(P_i, c_i, n_i)$, whose second term is the current payment (by the amortization identity, equal to the payment on $K_i$ at $c_i$ over the remaining term), the loss-aversion-scaled penalty, in \$/month, is term-dependent:
\begin{equation}\label{eq:abmpenalty}
\pi_{i,t} =
\begin{cases}
\lambda\,[\Delta_{i,t}]_{+} + [\Delta_{i,t}]_{-} & n_i > 240 \text{ (30-year book)},\\
\lambda\,\big[M^{\mathrm{new}}_{i,t} - A(K_i, c_i, n_i)\big]_{+} & n_i \le 240 \text{ (native short-term gate)},
\end{cases}
\end{equation}
% S:383-409; the 240-month switch is S:408
with $\lambda = 2.25$, $[x]_{+} = \max(x,0)$, and $[x]_{-} = \min(x,0)$. Payment increases are scaled by $\lambda$. On the 30-year book, decreases pass through unscaled as gains. The short-term gate instead floors them at zero, isolating the pure rate-lock cost. The frozen production headline draws 15-year voluntary CPR from the same-coupon 30-year surface, so the short-term branch is operative only in the Berger-run calibration stage (Section~\ref{sec:abm-danish}).
The effective transaction rate is $\tau_{i,t} = \mathrm{clip}\big(\phi_t + (\tau_i - 0.07),\,0.02,\,0.12\big)$. The agent's idiosyncratic draw enters as an offset around the 7\% base mean. It is then re-clipped after the macro friction $\phi_t$ shifts it. % S:241-243, S:459-461
Household $i$ moves in month $t$ iff
\begin{equation}\label{eq:abmmove}
\underbrace{\frac{M^{\mathrm{new}}_{i,t}}{y_i/12} \le 0.43}_{\text{DTI veto}}
\;\wedge\;
\underbrace{D_i > S\,\pi_{i,t} + \tau_{i,t} H_i}_{\text{cost--benefit}}
\;\wedge\;
\underbrace{\neg\big(v_t > 0.015 \,\wedge\, u_i < 0.20\big)}_{\text{wait-and-see freeze}},
\end{equation}
% Gates: DTI S:455-457, cost-benefit S:458-463, freeze S:464-467, conjunction S:468
where $D_i \sim \mathrm{Exponential}(\theta)$ is the mobility-desire draw in dollars, $S = 60$ months is the expected-stay horizon, and $v_t$ is the trailing six-month mortgage-rate change (decimal). The horizon converts the monthly penalty into a dollar stock. The remaining term, $u_i \sim U[0,1]$, is a fixed per-agent trait: the code's ``patience'' draw, consumed only by the freeze gate. It is not a time-preference parameter. % S:48, S:192, S:302-309
All inequalities are as coded: the DTI gate passes at exactly 43\%, the cost--benefit gate requires strict exceedance, and the freeze binds only strictly above 150 bp. At an evaluation point, the share of a pool passing \eqref{eq:abmmove} is read as that point's annualized CPR; the calibration anchor is an annual floor. The macro accounting layer then consumes it at $\mathrm{CPR}/12$ per month. % F:957

\paragraph{Parameter provenance.}
Table~\ref{tab:abmparams} collects every behavioral parameter. None is tuned to the \$764.7 billion benchmark. The mobility-desire scale $\theta$ is the one parameter re-derived under the floor rule, and only ever that one.

\begin{table}[H]
\centering
\footnotesize
\begin{threeparttable}
\caption{ABM behavioral parameters: value, units, distribution, and provenance.}
\label{tab:abmparams}
\begin{tabular}{@{}p{3.3cm}p{5.9cm}p{1.6cm}p{4.0cm}@{}}
\toprule
Parameter & Value / distribution & Units & Provenance \\
\midrule
Population & $N = 10{,}000$; engine seed 42 & households & author choice \\
Income $y_i$ & log-normal, $\sigma = 0.45$, median from FRED MEHOINUSA672N & \$/year & FRED-calibrated draw \\
Home value $H_i$ & log-normal, $\sigma = 0.35$, median from FRED MSPUS & \$ & FRED-calibrated draw \\
Origination LTV & 0.80, fixed & --- & author choice \\
Mobility desire $D_i$ & exponential, scale $\theta$ & \$ & floor recalibration (last row) \\
\quad $\theta$, production (frozen) & 43{,}883 & \$ & floor anchor, synthetic population \\
\quad $\theta$, cross-design recalibrated & 36{,}086 & \$ & floor anchor, real Freddie covariates (Section~\ref{sec:robustness-crossdesign}) \\
\quad $\theta$, external variant & 4{,}472 & \$ & \citepos{liebersohn2024} zero-gap 5.9\%/yr moving level, no floor (Section~\ref{sec:robustness-crossdesign}) \\
Expected-stay horizon $S$ & 60 & months & author calibration \\
Loss aversion $\lambda$ & 2.25 & --- & \citet{tversky1992} \\
DTI wall & 0.43, front-end & --- & 12 C.F.R. \S{} 1026.43(e)(2)(vi), behavioral anchor \\
Freeze share; trigger & 0.20 of agents; $v_t > 0.015$ over six months & ---; decimal & author calibration \\
Patience trait $u_i$ & uniform $[0,1]$, fixed per agent & --- & implementation device (freeze membership) \\
Transaction-cost draw $\tau_i$ & normal$(0.07,\,0.015)$, clipped $[0.02, 0.12]$ & share of $H_i$ & author calibration \\
Dynamic friction $\phi_t$ & $0.07$ base $+$ search penalty (listing shortfall vs.\ 2017--19 baseline, proportional, cap $+0.02$) $+$ sentiment penalty (5 bp per UMCSENT point below 85, cap $+0.015$); realized 8.23--9.96\% over the 42 QT months & share of $H_i$ & FRED ACTLISCOUUS, UMCSENT; mappings and base level are author calibrations \\
Curtailment & linear in real-income YoY growth: 0 at $-2\%$ to 1.2\% CPR at $+4\%$; $/12$ monthly, additive, macro layer only & annualized CPR & GSE daily-prepayment report estimates; FRED DSPIC96 (Appendix~\ref{sec:method-frictions}) \\
Turnover floor & binary search $\theta \in [10^{3}, 5\times10^{5}]$, at most 30 iterations, until CPR at an 8\% market rate under the U.S. rule, base friction, zero velocity lands in $[4\%, 5\%]$, evaluated on the dominant-weight 30-year SOMA cohort & annual CPR & 2023--24 discount-cohort turnover (Appendix~\ref{sec:method-frictions}) \\
\bottomrule
\end{tabular}
\begin{tablenotes}
\footnotesize
\item Committed artifacts carry the unrounded binary-search scales (43{,}882.8125; 36{,}085.9375; 4{,}472.05). The module default $\theta = 12{,}500$ is a placeholder overwritten by calibration at every run. Code pins: draws \texttt{abm\_lockin\_simulation.py} 296--309; constants 38--69; floor search 649--685; friction \texttt{fed\_mbs\_extension\_risk.py} 66--70, 294--334; curtailment 85--87, 141--155; variant scales \texttt{cross\_design\_test.py} 81--111 and \texttt{abm\_external\_gates.py} 88--89.
\end{tablenotes}
\end{threeparttable}
\end{table}

\paragraph{Order of operations.}
Within a simulated month, the three gates of \eqref{eq:abmmove} are evaluated on the surviving population and combined conjunctively. They are side-effect-free, so their order cannot change the result. The code evaluates them in vectorized form, in the order DTI, cost--benefit, freeze. % S:449-468
In cohort-path (burnout) mode, movers are removed permanently before the next month's evaluation. The turnover tail therefore sets a natural floor on the depleted pool. % S:470-518
Two named components are not in-month gates. The 4--5\% floor enters once, before simulation, through the desire-scale search of Table~\ref{tab:abmparams}. That search is re-run, for that parameter only, whenever a specification change alters the penalty or covariate structure facing every household. Because the anchor is evaluated at zero velocity, no freeze parameterization can move it (Appendix~\ref{sec:method-frictions}). % S:649-685, S:520-524
Curtailment never touches the household decision: it is an additive term in the macro layer's simulated roll-off. Nor does production re-solve the decision rule monthly. The CPR surface is precomputed on the rate $\times$ friction $\times$ velocity grid, with all gates applied at each node, then interpolated month-by-month. Appendix~\ref{sec:method-abm} discloses the freeze's step-versus-ramp semantics. % F:873-895, S:526-548

This paper earlier described the mobility-desire distribution as ``log-normal''. It is exponential, on a dollar scale. The correction is recorded in the replication package's run ledger.

\section{Path B Parameterization and Delinquency Transitions}\label{app:params}

\begin{table}[H]
\centering
\footnotesize
\begin{threeparttable}
\caption{Objects and notation for the two hazard specifications. Path A's macro coefficients are reported in Section~\ref{sec:patha}, and every Path B parameter with its provenance is collected in Table~\ref{tab:params}.}
\label{tab:hazard-notation}
\begin{tabular}{@{}lllc@{}}
\toprule
Symbol & Object & Units & Path \\
\midrule
\multicolumn{4}{@{}l}{\emph{Sample and exposure}}\\
$s$ & stratum (296 four-way cells) & -- & A \\
$(s,t)$ & unit of observation (stratum, month) & -- & A \\
$i$ & loan (75,000 sampled; 40,234 active) & -- & B \\
$(i,t)$ & unit of observation (loan, month) & -- & B \\
$y_{s,t}$ & prepaid unpaid principal balance & dollars & A \\
$n_{s,t}$ & UPB exposure (log offset) & dollars & A \\
\midrule
\multicolumn{4}{@{}l}{\emph{Link and baseline}}\\
$\alpha_{s}$ & stratum fixed effect (295 enter) & log-rate & A \\
$f(a_{s,t})$ & loan-age seasoning spline & log-rate & A \\
$h_{0}(a_{i,t})$ & PSA baseline hazard (Section~\ref{sec:pathb}) & monthly & B \\
$\underline{h}$ & baseline turnover floor & monthly SMM & B \\
$h^{\mathrm{vol}}$ & voluntary prepayment hazard, \eqref{eq:pathB} before the floor & monthly & B \\
$h^{\mathrm{prep}}_{i,t}$ & prepayment hazard, \eqref{eq:pathB} after the floor & monthly & B \\
$h^{\mathrm{def}}_{i,t}$ & competing default hazard, \eqref{eq:default} & monthly & B \\
\midrule
\multicolumn{4}{@{}l}{\emph{Covariates and coefficients}}\\
$g_{s,t}$ & rate gap & 100 bp units & A \\
$b_{s,t}$ & stratum burnout state & demeaned share & A \\
$\phi_t$ & dynamic friction index & index & A \\
$\beta_g,\ \beta_b,\ \beta_\phi$ & rate-gap, burnout, friction coeff. & log-rate & A \\
$g_{i,t}$ & rate gap (negative when locked in) & 100 bp units & B \\
$\tilde g_{i,t}=g_{i,t}/100$ & rate gap in decimal units & decimal & B \\
$\sigma_{i,t}=\max\{0,-\tilde g_{i,t}\}$ & rate stress & decimal & B \\
$z^{\mathrm{FICO}}_{i},\ z^{\mathrm{LTV}}_{i}$ & standardized static covariates & SD & B \\
$B_{s(i),t}\in[0,1]$ & cohort-level burnout state & cumulative share & B \\
$\beta_1,\ \beta_F,\ \beta_L,\ \beta_B$ & gap, FICO, LTV, burnout coeff. & log-hazard & B \\
$\gamma_r,\ \gamma_F,\ \gamma_L$ & default-hazard coefficients & log-hazard & B \\
\midrule
\multicolumn{4}{@{}l}{\emph{Convention and transport parameters}}\\
$\mu$ & moving-share bracket (production $\mu = 1$) & share & B \\
$\omega$ & floor-form mixture weight (production $\omega = 0$) & share & B \\
$\underline{h}_t = 4\%\,(1+\kappa z_t)$ & within-run cyclical floor variant & monthly SMM & B \\
$\eta = S_{\mathrm{pre}}^{\xi}$ & survival-selection attenuation & factor & B \\
$S_{\mathrm{pre}}$ & pre-window survival share & share & B \\
$\xi$ & gamma-frailty parameter (not estimable here) & -- & B \\
\bottomrule
\end{tabular}
\begin{tablenotes}[flushleft]\scriptsize
\item Case and subscript separate the two paths throughout, and no coefficient symbol is reused across them. That guarantee is about coefficients: the convention and transport parameters in the last block bracket reported results rather than entering either fit, and none of them shares a glyph with the stratum index $s$ or with each other. Path A's burnout coefficient $\beta_b$ (on $b_{s,t}$, estimated on the panel) and Path B's $\beta_B$ (on the cohort state $B_{s(i),t}$, a calibrated prior) are distinct objects carrying distinct symbols, as are Path A's friction coefficient $\beta_\phi$ (on the friction index $\phi_t$) and Path B's FICO coefficient $\beta_F$. The rate-gap elasticities $\beta_g$ (Path A, log-rate, standardized) and $\beta_1$ (Path B, log-hazard, per 100 bp) are on different scales; the paper compares them on sign, not magnitude (Section~\ref{sec:hazard-interp}).
\end{tablenotes}
\end{threeparttable}
\end{table}

{\footnotesize
\begin{longtable}{@{}p{1.5cm}p{6.2cm}p{6.2cm}@{}}
\caption{Reading key for words this paper reuses across designs. Each row gives the sense an unqualified occurrence carries and the qualifiers that mark every other sense in use. Symbols are in Table~\ref{tab:hazard-notation}; no row states a result, and every distinction below is made where it is used.}
\label{tab:terminology}\\
\toprule
Word & Unqualified, it denotes & Marked senses it must not be read as \\
\midrule
\endfirsthead
\toprule
Word & Unqualified, it denotes & Marked senses it must not be read as \\
\midrule
\endhead
stratum & Path A's four-way panel cell, vintage $\times$ coupon $\times$ FICO $\times$ LTV: 296 observed, 295 fixed effects, and the cluster unit that carries the floor reads (31 occupied off window, 226 in window) & the draw stratum: the same key with the LTV axis dropped, 130 of them and 25.8 effective, the unit of Path B's loan bootstrap. Both are resampled under the name ``stratum-cluster''; the path decides which \\
cell & nothing on its own --- the partition is always named, and the counts do not interchange & a panel stratum (296 four-way; 130 three-way for the draw); a stratum-month observation; the standardization cell of Section~\ref{sec:patha}, the panel key with the coupon axis removed (22 cells, 29 with 12-month age bands, 18 in common); a point of the 21-cell calibration box (Table~\ref{tab:floorband}: seven floors $\times$ three elasticities); a point of the twenty-seven-cell off-window floor $\times$ elasticity-band grid (run \texttt{oos\_identification}); a joint agency $\times$ vintage cell of the SOMA book; a printed table cell. ``Every cell positive'' is claimed once over book subgroups and once over the off-window grid, and those are two different claims \\
band & the imported rate-response range, 5.5\%--7.7\% with 6.5\% adopted --- a calibration interval, not sampling error & the 4--5\% turnover band, which is a floor \emph{level} range; the clean band, the admissible off-window floor reads; the seed band, the $\pm 1$~SD band and the moving-block bands, which are sampling dispersion; the corroborating band, a pre-committed acceptance range; age bands and single-band vintages, which are bucket edges \\
anchor & in the hazard sections, the floor level a read is taken at: in-window versus off-window, the mid-grid anchor being the middle of the three depth cuts fixed before the reads & the transplant construction of Section~\ref{sec:abm-danish}, U.S.-intercept versus Danish-level; an outside estimate quoted beside the band; the ABM's mobility-scale anchor, set by the floor search; anchor-gated, meaning checked against the committed SOMA marginals \\
basis & the accounting basis: shared, standalone, full-book, composed (Table~\ref{tab:bases}). ``Basis-invariant'' means invariant across those four and across nothing else & the coupon convention behind a CPR back-out --- ABM basis, hazard basis, note-rate basis --- which does move the printed levels; the weighting behind a share (count, origination balance, surviving balance); basis points; the seasoning spline's power basis \\
floor & the baseline turnover floor $\underline{h}$, a level in annual CPR & the floor \emph{form}, the rule combining it with the voluntary hazard, hard maximum versus additive competing risks --- an axis, not a level; the floor \emph{read}, the measurement of that level off discount cohorts; the scheduled-amortization floor the back-out's non-negativity clip pins against \\
form & the floor form above; ``form-conditional'' is always this axis & the rendering of an estimator, microsimulation form versus cohort-hazard form of the same framework; functional form of a hazard in general; decimal form and CPR form, where only a restatement of one object is meant \\
draw & the stratified 75{,}000-loan sample Path B simulates & a bootstrap or permutation replicate, always plural and counted; a random variate inside a run, such as the patience draw or the default draw; a seed-level realization of the ABM, the frozen draw; draw seeds, which are simulation seeds and not resamples of that sample \\
null & the $\beta_1 = 0$ rate-inelastic leg: a paired simulation run scored in dollars & a null hypothesis --- the $H_0$ tests of Appendix~\ref{app:patha-detail}, and the mean-only null of this paper's $R^2$ convention. A leg that clears a threshold is not a non-rejection \\
\bottomrule
\end{longtable}
}

\noindent{\footnotesize \emph{Notes to Table~\ref{tab:terminology}.} Rows are ordered by how badly a conflation would mislead, most damaging first. ``Cluster'' is deliberately absent. It carries one sense here, the resampling unit of a cluster-robust bootstrap, and what varies is the unit, not the word --- the four-way panel stratum, the three-way draw stratum, the six calendar months of the off-window leg, and the 137 cells where those two intersect. The stratum and cell rows fix each.\par}

Table~\ref{tab:params} collects every Path B parameter. All values are fixed a priori in the configuration, and none is estimated on, or tuned to, the \$764.7 billion benchmark. One parameter is nonetheless measured inside the evaluation window, and benchmark-independence should not be read as window-independence. The 4\% baseline turnover floor is anchored to turnover observed on discount cohorts during 2023--2024 (Appendix~\ref{sec:method-frictions} and Section~\ref{sec:pathb}), so it is in-sample with respect to the QT window even though it never touches the benchmark dollar figure. That floor sets much of the aggregate level, binding in 36\% of loan-months. Section~\ref{sec:robustness-floor} sweeps it jointly with the elasticity band: the aggregate level and the lock-in marginal both move materially with the floor, within the bounded interval that section reports. The marginal's sign does not move, but under this window's rate configuration it could not (Section~\ref{sec:identification}). The covariate and default-hazard coefficients are round calibration priors. Both prepayment covariates carry signs opposite the canonical credit-risk direction (higher FICO $\Rightarrow$ better refinancing access $\Rightarrow$ faster prepayment; higher LTV $\Rightarrow$ less equity $\Rightarrow$ slower prepayment). In the deeply out-of-the-money 2022--2025 window the refinancing channel is shut and prepayment is turnover-dominated, so credit covariates proxy mobility, and high-LTV loans are disproportionately recent purchases that churn. A Path-A-style estimation on the same stratum-month panel yields point estimates consistent with the priors' direction; taken at face value, they are roughly double the priors' magnitudes. That estimation is a Poisson PML with vintage-by-coupon fixed effects, so the credit buckets are identified from cross-bucket variation, and the bucket-mean covariates are z-scored to the microsim pool's own moments. But the intervals reject neither zero nor the opposite signs: $\hat{\beta}_F = -0.39$ (four-way stratum-cluster bootstrap 95\% CI $[-1.51, +0.65]$) and $\hat{\beta}_L = +0.21$ ($[-0.59, +1.00]$). The estimation does not settle the sign question; it fails to contradict the priors while quantifying their sensitivity. The sensitivity brackets are correspondingly modest: a $\beta_F = \beta_L = 0$ rerun moves the recovery from 107.0\% to 108.3\%, and the panel-estimated pair moves it to 101.3\%. Those are in-window figures. At the off-window floor the same endpoints span 8.2 points and `modest' does not carry over to the level, while the marginal's covariate sensitivity shrinks to $+0.1$ points. Benchmark-independence is unaffected, since the estimation touches only the Freddie panel. The $\beta_B = 0$ rerun of Section~\ref{sec:pathb} quantifies the burnout channel the same way.


Delinquent loans transition monthly through the matrix of Table~\ref{tab:delinq}. I estimate it from consecutive stratum-month \emph{modal} servicer states in the Freddie cohort panel, exposure-UPB weighted. I commit it as an artifact in the replication package. The default-hazard draw and the transition matrix divide the labor. The draw moves loans \emph{into} the pipeline (Current $\to$ D30) and terminates already-delinquent loans it hits. The matrix governs month-to-month movement \emph{within} the pipeline. It runs after the draw each month, so a loan entering D30 via the draw can also take a matrix transition in its entry month. The \$0.39 billion extreme-matrix test below bounds that overlap along with everything else. Only the delinquency rows are used in simulation. Terminal states are absorbing by construction. The forbearance state is unpopulated in the estimation panel and never entered in production runs.
I publish the full crosstab. D30 row (21 transitions, \$58.3M row exposure): 13 cures carrying 95.6\% of the exposure weight, two D30 continuations (0.8\%), five D60 progressions (3.1\%), and one Defaulted transition (0.4\%). The printed 0.004 is the exposure share of that single event; its count share would be $1/21 = 0.048$. D60 row (9 transitions, \$3.1M): six cures carrying 56.8\% and three D90+ progressions carrying 43.2\%, against count shares of $6/9$ and $3/9$. That divergence between count and exposure shares means a few larger-balance stratum-months drive the matrix. That is what an exposure-weighted estimator is designed to do. D90+ row (38 transitions, \$76.2M): seven cures (30.3\%), one D30 (0.8\%), thirty self-continuations (68.9\%), and zero Defaulted events. The ``--'' cells are observed zeros in a thin panel, not structural zeros. That is why serious delinquency shows no direct default exit while D30 shows one. Probabilities are printed to three decimals from these exposure shares, and rows may not sum to exactly one. The matrix is accordingly reported for completeness, not precision: with cells this thin it describes the estimation panel and does not estimate a transition law. The paper's estimates rely on the materiality bound below, not on the matrix's point values.

The thinness is immaterial for the paper's estimates. Replacing the entire matrix with either extreme --- all delinquent loans cure next month, or all default next month --- moves trapped liquidity by \$0.39 billion (0.05\% of the estimate). The matrix routes the 482 loans delinquent at the window start from the first month, and given the cure-heavy rows it mostly cures them. Thereafter the month-end delinquency stock never exceeds 154 of the 40,234 active loans, and delinquent UPB never exceeds 1.1\% of the pool. Both extreme-matrix runs include the initial 482, so the bound covers them.

\begin{table}[H]
\centering
\small
\begin{threeparttable}
\caption{Exposure-weighted delinquency transition matrix (modal stratum-month states).}
\label{tab:delinq}
\begin{tabular}{@{}lrrrrr@{}}
\toprule
From $\backslash$ To & Current & D30 & D60 & D90+ & Defaulted \\
\midrule
D30 & 0.956 & 0.008 & 0.031 & -- & 0.004 \\
D60 & 0.568 & -- & -- & 0.432 & -- \\
D90+ & 0.303 & 0.008 & -- & 0.689 & -- \\
\bottomrule
\end{tabular}
\begin{tablenotes}[flushleft]\footnotesize
\item Probabilities are exposure-UPB shares of thin modal-state transition counts (D30 row: 21 transitions; D60: 9; D90+: 38; full crosstab in the text and replication artifact). ``--'' cells are observed zeros in a thin panel, not structural zeros; rows may not sum to exactly one at printed precision. Reported for completeness, not precision: the paper's estimates rely on the \$0.39 billion extreme-matrix materiality bound in the text, not on these point values.
\end{tablenotes}
\end{threeparttable}
\end{table}
Sampling design: the draw is not a probability sample of the SOMA book, and the allocation rule is why. A pool built at four times the target size takes an equal allocation from each quarterly Freddie origination file. That pool is then drawn down to 75,000 loans proportionally within stratum, and every retained loan carries weight one (\texttt{hazard/loan\_sample.py}). The draw's stratum is the three-way vintage $\times$ coupon $\times$ FICO cell --- the Path A estimation stratum with the LTV axis dropped. That is why the draw carries 130 of them against the panel's 296. Nothing in that chain reweights toward the book. Each origination year from 2017 through 2021 therefore takes 20.0\% of the draw by count, so 2017--19 is 60.0\% of it against 6.0\% of book face. But count is not the weighting the aggregation applies. The aggregation weights by surviving balance at the window open, and on that basis the picture is materially different: 2017--19 is 27.0\%, 2020 30.6\% and 2021 42.3\%, against book face shares of 6.0\%, 16.3\% and 43.9\%. So the 2021 cohort is close to matched. The 2017--19 overweight is a factor of about four and a half, not the factor of ten the count basis suggests. Origination balance gives 55.5\%, 22.2\% and 22.3\%, but it compares originations to current face and is the wrong basis for this comparison. What no weighting repairs is out-of-support mass. The draw spans 2017 through 2021 only. So pre-2017 (10.6\% of book face) and the 2022 cohort (23.1\%) carry no support at all: 33.7\% of face, which must be netted against the Ginnie share on the book's own joint cells, not added to it. The draw's 3.86\% unweighted mean note rate also sits above the book's 2.49\% face-weighted pass-through coupon. Both tilts move the rate-gap distribution the floor takes its maximum against under the production floor form, and this paper signs neither. Section~\ref{sec:robustness-hybrid}'s coupon reweight rescales the aggregate level and leaves the marginal unchanged, so it does not test a composition effect running through the floor's bind share.

Attrition accounting: of the 75,000 sampled loans, 34,734 prepaid and 32 defaulted before the QT window opens. That leaves 40,234 active at the June 2022 start --- 39,752 current, 482 delinquent. The production U.S. run evaluates 1,683,124 loan-months over the 42-month window, a mean of 40,074 active loans per month. The baseline turnover floor $\underline{h}$ binds in 36\% of those loan-months (Section~\ref{sec:pathb}).

The $\delta$-to-$\beta_1$ conversion of \eqref{eq:beta1} is insensitive to the turnover proxy behind it. In the small-$P_q$ limit the transform reduces to $\beta_1 = -\ln(1-\delta) = 0.0672$ at the central $\delta$. Over $P_q \in (0, 0.12]$---from arbitrarily low turnover to double the adopted proxy---the central $\beta_1$ ranges only from $0.0672$ to $0.0701$. Even a large error in $P_q$ therefore moves the conversion by a few percent.

\paragraph{Production step order, clamp, and parameters.}
Each month, both hazards are evaluated for every active loan. In the production configuration prepayment is fractional: each active loan's balance is drained deterministically by $h^{\mathrm{prep}}$ (a per-loan uniform draw for prepayment is a non-production discrete branch). Default instead turns on a uniform draw. Hits enter the delinquency pipeline --- current loans to D30, already-delinquent loans to default --- where they transition through the Markov delinquency-resolution matrix of Appendix~\ref{app:params}. A proportional rescaling guards the case $h^{\mathrm{prep}} + h^{\mathrm{def}} > 1$. It is a safety clamp, not a competing-risks construction. In production it never binds: across all 1,683,124 evaluated loan-months per regime the maximum hazard sum is 0.016, two orders of magnitude below the threshold. At that level the exact exponential cause-allocation construction $1-\exp(-(\lambda_p+\lambda_d))$ with shares $\lambda_p/(\lambda_p+\lambda_d)$ differs from mine only at second order in the hazards. In the taxonomy of standard competing-risks formulations \citep{putter2007}, Path B is a discrete-time multi-state model parameterized entirely through cause-specific hazards. Equations~\eqref{eq:pathB} and \eqref{eq:default} give monthly event probabilities conditional on the loan remaining active --- the estimand of discrete-time cause-specific regression, as implemented in open-source packages such as PyDTS \citep{meir2025}. The delinquency pipeline of Appendix~\ref{app:params} extends that two-cause structure to a multi-state process with cures. No subdistribution (Fine--Gray) hazard \citep{fine1999} appears in the specification, and none is needed. That device places covariate effects directly on cumulative incidence functions. The simulation instead obtains every cumulative-incidence functional --- prepaid UPB, the delinquency stock, trapped liquidity --- by forward composition of the cause-specific hazards. One caveat qualifies the mapping. In production, prepayment is a fractional balance-flow intensity, not a per-loan absorbing transition, so the cause-specific reading is exact for default and holds at the balance-aggregate level for prepayment. Prepaid balances settle to simulated SOMA cash flow in the month they occur; no settlement delay is modeled on this path.

\begin{table}[H]
\centering
\footnotesize
\setlength{\tabcolsep}{4pt}
\begin{threeparttable}
\caption{Path B parameters.}
\label{tab:params}
\begin{tabular}{@{}llll@{}}
\toprule
Parameter & Value & Units & Provenance \\
\midrule
$\beta_1$ (central) & $0.069$ & per 100 bp gap & \eqref{eq:beta1} at $\delta = 0.065$ \citep{liebersohn2024} \\
$\delta$ band & $0.055$--$0.077$ & quarterly decline & \citet{liebersohn2024} \\
$P_q$ & $0.06$ & quarterly slot & turnover proxy (conversion auxiliary; not a measured quarterly mobility rate) \\
$\beta_F$ (FICO, z) & $-0.15$ & per SD & calibration prior; panel estimate $-0.39$ (see text) \\
$\beta_L$ (LTV, z) & $+0.10$ & per SD & calibration prior; panel estimate $+0.21$ (see text) \\
$\beta_B$ (burnout) & $-0.5$ & per unit of $B \in [0,1]$ & calibration prior \\
$\underline{h}$ (floor) & $0.0034$ & monthly SMM & 4\% annual CPR (baseline turnover) \\
PSA speed & 100 & -- & industry seasoning convention \\
$h^{\mathrm{def}}_0$ & $3\times10^{-4}$ & monthly & calibration prior \\
$\gamma_r$ (rate stress) & $2.0$ & per decimal unit & calibration prior ($\approx$2\% per 100 bp) \\
$\gamma_F$ (FICO, z) & $-0.20$ & per SD & calibration prior \\
$\gamma_L$ (LTV, z) & $+0.25$ & per SD & calibration prior \\
\bottomrule
\end{tabular}
\end{threeparttable}
\end{table}

\paragraph{Cohort-level burnout.}
Burnout is deliberately cohort-level, not per-loan: a fractional burnout measure is not well-defined for a one-time, absorbing transition (Section~\ref{sec:abm-results}). Each stratum instead tracks its cumulative prepaid UPB as a share of its original UPB---the $B \in [0,1]$ state defined under \eqref{eq:pathB}---after all prepayment transitions in a given month. That state is broadcast to all of that stratum's remaining active loans before the next month's hazard is evaluated, and enters \eqref{eq:pathB} with the calibrated prior coefficient $\beta_B = -0.5$ (Appendix~\ref{app:params}). Path A's separately estimated burnout coefficient, statistically indistinguishable from zero (Section~\ref{sec:patha}), plays no role here. Because the simulated burnout state is endogenous, I test whether the channel is load-bearing with a $\beta_B = 0$ rerun on the production seed: trapped liquidity moves by \$6.58 billion ($-0.86$pp of benchmark, to 106.17\%), with the peak lag unchanged at $-3$. Two qualifications travel with that figure. It is measured at the in-window 4.0\% floor; at the 4.991\% off-window floor that the paper headlines, the same ablation is worth only \$3.48 billion ($-0.45$pp), because the floor's bind share rises from 36.3\% to 68.8\% of evaluated loan-months and the hard maximum makes burnout inert wherever it binds. And because the burnout state is endogenous to realized prepayment, the ablation moves the $\beta_1 = 0$ null as well as the central leg, so the identified marginal is not burnout-invariant: switching the coefficient off raises it from $+9.20$ to $+9.94$ points at the production floor and from $+5.57$ to $+6.51$ at the off-window floor. Both moves sit inside the $\pm 1$-point materiality convention, the second only barely, so I report them without resting anything on them. The FICO and LTV covariates receive the same treatment. A $\beta_F = \beta_L = 0$ rerun moves the recovery to 108.3\% ($+1.3$pp), and a rerun under coefficients estimated Path-A-style from the panel moves it to 101.3\% ($-5.7$pp), both at the in-window 4.0\% floor. At the 4.991\% off-window floor the same endpoints span more, not less: 102.5\% against 94.3\% standalone, an 8.2-point swing. The marginal they bracket nevertheless moves by only $+0.1$ points, against $+1.9$ in-window, because both legs shift together where the floor censors more. Appendix~\ref{app:params} reports the estimation, whose imprecise point estimates are consistent with the priors' turnover-regime signs without rejecting alternatives. Production retains the ex-ante priors---the estimates are too imprecise to license replacing a pre-specified calibration mid-stream---and I flag the covariate block as the specification's largest single sensitivity: the two endpoints correspond to 99.2\% and 92.2\% on the shared accounting basis (the common netting makes the conversion exact), a swing comparable in size to the lock-in marginal itself.

\paragraph{The time-varying Ginnie overlay.}
The static Ginnie bound is now extended to a time-varying overlay built from the published series the bound's two snapshots came from. I extract the full monthly agency CPR series (June 2017 through November 2025) from the December 2025 Global Markets Analysis Report, and score the 20.4\% Ginnie face share by the observed Ginnie speed in \emph{both} Path B legs. That report's aggregation was revised in that issue to weighted-average unpaid principal balance, one consistent dollar-weighted methodology across the whole window. In the raw construction the overlay moves the shared-basis central from 97.9\% to 89.3\% and the null from 88.7\% to 81.9\%, both legs at the in-sample calibration point. It confines the identified marginal to the conventional share: $+7.3$ points, positive under every variant tested (total-CPR, voluntary-only, and placebo), scaling with the conventional face share almost exactly ($0.797$ against $0.796$). Composed with the off-window headline calibration, the overlay pair is 84.0\%/79.5\% central/null on the shared basis with a $+4.4$-point marginal. That overlay marginal is pure conventional-share scaling ($0.797\times$) --- a change of estimand, a re-scoring onto the conventional sub-book, identical across the primary, voluntary-only, and placebo series. It is not a correction to the headline. The headline marginal therefore carries the conventional-extrapolation convention wherever quoted, with $+4.4$ points its overlay-scored member. The raw level shift overstates the Ginnie-specific effect, though. A placebo that scores the same 20.4\% share by the observed GSE-mean series moves the level by $+\$27.8$ billion on its own: the model's level calibration plus the full-universe-versus-SOMA composition difference, common to any observed-series overlay. The Ginnie-specific differential contribution is therefore $+\$38.5$ billion, 5.0\% of benchmark. That sits inside the static \$20--\$47 billion bound above. It is also consistent with the window-mean Ginnie--Freddie differential of 2.1 points the series itself shows, against the bound's 1.2--2.8-point snapshots. The placebo and Ginnie-specific components, $\$27.8 + \$38.5 = \$66.3$ billion, reconcile with the raw 8.6-point shared-basis shift on the \$764.7 billion benchmark, $\approx\$65.8$ billion, to within rounding of the 0.1-point recovery shares. The full-universe series if anything overstates the correction for SOMA's seasoned, low-coupon Ginnie holdings, the same signed direction as the static bound. Structural Ginnie inclusion remains out of scope for the modeling reason above: the overlay treats the buyout-inclusive Ginnie speed as data, which is what runoff accounting requires of a passthrough holder. This overlay and the age-standardized floor read of Section~\ref{sec:robustness-floor} cut the same way, and the joint cell has now been run under a pre-committed landing rule (run \texttt{b5\_joint\_cell}). Composed, they return $+2.9$ points ($+\$22.4$ billion), within $0.1$ points of the proportional projection, and that measured point is the assembly's lower member. Pure conventional-share scaling is what a gap-inert Ginnie share gives. Scoring that share at a response instead (run \texttt{ginnie\_overlay\_attenuated}) prices it at a fixed proportion of the conventional share's per-unit-face wedge (transplanted, not estimated on Ginnie collateral), with the wedge fixed by requiring that a full-strength Ginnie response return the un-overlaid marginal. What the series measures is a level, and a voluntary-only one: across the window's 42 months Ginnie's voluntary speed runs $1.35$ points above Freddie's, $1.228$ times as fast, against the 2.1-point total-CPR differential above, and read off the same full-universe series with the same overstatement. An elasticity is a slope, and nothing here maps a level onto one, so I bracket the mapping rather than choose it. The four are: zero (the committed overlay); $81.4$\% of the conventional response, the reciprocal of that speed ratio, on which a book running $1.228$ times faster carries a proportionally smaller semi-elasticity; $100$\% (no differential); and $122.8$\% (proportional-hazard preservation, which is amplification, not attenuation). In-sample the marginal runs $+7.3$, $+8.9$, $+9.2$ and $+9.6$ points across those four; composed with the off-window calibration, $+4.4$, $+5.4$, $+5.6$ and $+5.8$. That the scaling rises above the committed $0.797\times$ at every positive response is arithmetic, not evidence---a share contributing exactly zero must contribute more at any positive response---and I rest nothing on it. The bracket, not any cell in it, is the result: none of the four is a measured Ginnie marginal, and the overlay-scored $+4.4$ stands unrevised, scoped by the bracket rather than corrected by it.

\paragraph{Danish pipeline mechanisms.}
Two Danish counterfactual mechanisms exist in this framework's pipeline. Only one is live in production. The first is a \emph{rate-gap swap}. Holding the baseline hazard, rate-gap coefficient, covariates, and random draw logic fixed, the Danish leg expresses the market-value buyback in rate-equivalent units, so the buyback discount offsets the locked-in spread. The swap sets $g_{\mathrm{DK}} = g_{\mathrm{US}} + \Delta y_{\mathrm{buyback}} \approx 0$ whenever the loan is out of the money. When rates have fallen and the buyback price is capped at par, $g_{\mathrm{DK}} = g_{\mathrm{US}}$. At the representative 3.0\% coupon / 6.8\% market state the U.S. gap is $-380$ bp while the Danish effective gap is zero. The U.S. hazard is therefore suppressed by the full locked-in spread. The Danish hazard resets toward its baseline seasoning curve. This swap defines the \emph{mechanism-substitution} variant, the specification behind the superseded 47.1\% Danish CPR of Section~\ref{sec:abm-danish}. It is \emph{not} the production mechanism. The production Danish leg is the Berger-recalibrated specification adopted in Section~\ref{sec:abm-danish}: it replaces the rate-gap swap with the imported flat moving elasticity, so the production Danish hazard does not respond to the buyback discount at all. The swap machinery remains in the codebase, and its PV-versus-par classification is still computed each month, but no production quantity consumes its output. That is why the discount-rate diagnostic of Appendix~\ref{sec:method-abm} moves the Danish leg by exactly \$0.00 billion. Table~\ref{tab:discount} reports that inertness. Shifting the buyback discount rate down by 25 to 150 bp---through the primary--secondary spread range---flips thousands of monthly present-value-versus-par classifications while moving the Danish CPR path, roll-off path, and trapped-liquidity dollar by exactly zero on every scoring basis. Two limits apply to that result. I state both because together they retract the immunity I previously claimed from it. First, the diagnostic executed on the Danish-level bracketing leg, not on the production U.S.-intercept leg: its own artifact records a mean Danish CPR of 3.39\%, the bracketing anchor's speed. Under the production anchor the same independence does hold, but by construction, not by measurement, because that branch replaces the rate gap with zeros before the hazard is evaluated. Second, and more important, the result is about the hazard, not about the accounting. The Danish leg credits every retired balance at par. Its dynamic-balance loop books roll-off as balance times drain, with no price term anywhere in it. Netting the roughly \$199 billion scheduled component from the leg's \$669.3 billion of Danish roll-off leaves on the order of \$470 billion of buyback-triggered retirement credited at par, against a committed proxy discount grid of 32 to 38\%. A market-value credit would reduce the Danish leg's cash receipts by an amount larger than the $+\$61.2$ billion gap itself, and in the direction that shrinks it. The pipeline does not distinguish a cash haircut from a balance adjustment. The incidence depends on which the buyback is. So a pre-committed bracket prices both readings from the committed pair (run \texttt{buyback\_credit\_bracket}; gates and the interpretive rule fixed in the script header before execution). Under the balance-adjustment reading face retirement is unchanged and the gap stands at $+\$61.2$ billion. Under the cash-haircut reading receipts fall by the discount on the leg's \$470.7 billion of early face and the gap reverses. The committed grid was asserted, not derived. Run \texttt{buyback\_discount\_rederived} prices each month's retired face at its own prepayment-consistent present value---the SOMA book aged to that month at the 2.49\% weighted-average coupon, prepaying at that month's Danish CPR, discounted at that month's market rate. That returns a face-weighted mean of 23.8\%, a \$112.2 billion haircut, and a cash gap of $-\$51.0$ billion. The price and the face come from two populations, the SOMA book's own vintage tabulation and the microsimulated leg's roll-off: the seam the committed early-face netting already carries, inherited here rather than repaired. That is a point, not a range: the committed 32--38\% span was a range only by virtue of a four-point asserted grid. The monthly discounts do range, from 18.4\% to 28.1\%, and none of the forty-two months reaches the committed grid's floor. The committed grid is consistent with a no-prepayment annuity priced at a representative state: the same primitive returns 35.3\% at zero prepayment on a 3.0\% coupon against a 6.8\% market rate, inside the committed span, where the leg itself prepays at 5.61\% against a window averaging 6.61\%. That identification is plausible rather than exactly reproduced: a band check on the primitive, not a target. The re-derivation shrinks the reversal without retiring it. The threshold it clears is arithmetic on the committed chain, fixed before the run rather than produced by it: the verdict turns at a 13.0\% discount, below every month's priced value. The two readings are not symmetric answers to one question. The benchmark is a face-value object. Section~\ref{sec:method-benchmark} constructs it from the SOMA current face value series against the par-denominated caps. So the balance-adjustment reading is the benchmark-consistent one, and on the benchmark's own denomination the gap stands. The cash-haircut reading answers a different question, reserves actually drained, and its reversal has a name. The $1-P$ wedge on bought-back face is exactly the par windfall a moving household pays the bondholder under the U.S. rule and keeps under the Danish one. So the cash-flow cost is the disappearance of a household-to-bondholder transfer, not an added resource cost. The gap's sign is therefore incidence-conditional, the channel is priced, and the discount-rate diagnostic still should not be read as immunity to the pricing critique.

\paragraph{The transform's superseded variant and the concave-transform test.}
Two records travel with the transform of \eqref{eq:beta1}. The superseded $0.0693$ is a distinct object: the continuous-hazard variant $-\ln\!\big[\ln(1-(1-\delta)P_q)/\ln(1-P_q)\big] = 0.0693$, unrounded $0.069258$. The simulation does not use it. The two forms differ by $0.0007$ and both round to $0.069$. Only \eqref{eq:beta1} reproduces the reported results.

The direct test of the log-linear extrapolation is a piecewise-concave transform---full marginal effect over the first 200 bp of lock-in, half beyond. It moves the recovery from 107.0\% to 105.7\% (\$808.4 billion) and slightly improves the contemporaneous path correlation ($+0.281$ versus $+0.190$, peak lag unchanged at $-3$). The log-linear extrapolation is not load-bearing for the level, and is for the marginal. But the asymmetry is one of proportion, not of magnitude, and the same absolute $\pm 1$-point convention cannot support both readings. A pre-committed rerun of the paired legs under the concave transform puts the concave marginal at $+7.87$ points against the production $+9.20$ at the in-sample floor, and at $+5.06$ against $+5.57$ at the off-window point floor. At the in-sample floor the concave null is bit-identical to the production null---\$748.185 billion on both legs, because $\beta_1 = 0$ leaves nothing for the transform to act on. So the level move and the marginal move are numerically the \emph{same} 1.32 points. That 1.32 points is 1.2\% of a 107.0\% recovery and 14\% of a $+9.20$-point marginal. It is that relative reading, not the absolute one, on which the transform is immaterial to the level and material to the marginal. On the absolute $\pm 1$-point materiality threshold the floor-form test of Section~\ref{sec:robustness-floor} applies, and both moves clear it. The transform dimension travels with the marginal's operative uncertainty in Table~\ref{tab:uncertainty}.

\paragraph{Sampling schemes for Path B's level and marginal.}
The Liebersohn--Rothstein band of Section~\ref{sec:pathb} is a calibration interval, not sampling error. The sampling dimension is quantified separately, by a stratified loan-level bootstrap specified ex ante in its script header: 200 replicates resampling the 75,000-loan draw with replacement within strata, central and null legs paired per replicate, engine seed held at production. That bootstrap puts the central recovery's 95\% interval at $[107.0, 107.1]\%$ under the standalone scorer (replicate standard deviation 0.03 points), the null's at $[97.8, 97.9]\%$, and the paired lock-in marginal's at $[+9.17, +9.23]$ points (standard deviation 0.016). Re-run under the cluster scheme at this same floor, the paired marginal reads $[+8.27, +10.19]$ points (sd 0.52; run \texttt{bootstrap\_pathb\_cluster}), an understatement of 30.8$\times$ in-window against 37$\times$ off-window. Those intervals are narrow because the scheme holds the design fixed. It resamples loans within each stratum and preserves every stratum's count, so between-stratum composition variance is never drawn. Drawing the 130 three-way draw strata themselves with replacement---at the off-window floor the paper headlines---widens the marginal's interval to $[+4.63, +6.92]$ points with a standard deviation of $0.60$ points (run \texttt{bootstrap\_pathb\_cluster}, 200 replicates). That is a width of $2.29$ points against the within-stratum scheme's $0.062$, so the committed scheme understated loan-sample uncertainty by roughly a factor of $37$. Two properties of that interval are not corrected for and should be stated. Drawing strata with replacement makes the loan count itself random: it runs from 54,134 to 101,404 across replicates against a fixed-design 75,000. So every composition weight is a ratio of random totals, not of fixed ones, and the percentile interval carries no correction for the resulting ratio bias. A fixed-$n$ variant that rescales weights to the design total would settle how much of the width is that bias. It is not run here. The committed point marginal falls inside that interval. Two limits apply to the wider figure. The first is that the strata are very unequal, so the effective cluster count is $25.8$, not $130$ (the largest single stratum is $9.77\%$ of sample balance), and a percentile interval on $26$ effective clusters is itself noisy. The second is that the resampled loan count varies from roughly $54{,}000$ to $101{,}000$. That variation is harmless here only because the engine renormalises the pool to Fed holdings each month, so the level is set exogenously and only composition varies.

\paragraph{Curtailment-layer demonstration runs.} On the accounting identity of Section~\ref{sec:robustness-hybrid}: a $\pm 25\%$-amplitude seasonal profile and an income-elasticity regime split across 2022--23 and 2024--25 move the netted total to \$70.06 and \$68.09 billion and the wedge to 9.16 and 8.90 points, the marginal unchanged to below $10^{-13}$ points on both legs and the CPR paths bit-identical; a deterministic two-servicer mix (fixed 0.60/0.40 shares, front-loaded and flat group factors) moves the wedge to 9.30 points (\$71.12 billion netted) with the marginal again unchanged, since any servicer-heterogeneous pattern aggregating to one monthly rate is a special case of the same identity.

\section{Path A Estimation Detail and Timing Diagnostics}\label{app:patha-detail}

This appendix collects the estimation detail of Path A (Section~\ref{sec:patha}), the direct input-timing diagnosis, the Danish discount-rate diagnostic, and the monthly-timing exhibits of Section~\ref{sec:identification}.

The second hazard path aggregates Freddie Mac loan-level performance data to the stratum-month level. Strata $s$ are the 296 observed four-way cells of vintage $\times$ coupon $\times$ FICO bucket $\times$ LTV bucket. The unit of observation is the stratum-month pair $(s,t)$. The estimation configuration sets the bucket edges: FICO below 680, 680--740, and 740 and above; LTV at or below versus above 80; coupons in 0.5-percentage-point steps; vintages 2017--2021. Geography is not a stratum axis, and no geographic term enters \eqref{eq:pathA}. The model is a discrete-time hazard with a log-exposure offset,
\begin{equation}\label{eq:pathA}
\log \mathbb{E}\!\left[y_{s,t} \mid x_{s,t}\right] = \log n_{s,t} + \alpha_{s} + f\!\left(a_{s,t}\right) + \beta_g\, g_{s,t} + \beta_b\, b_{s,t} + \beta_\phi\, \phi_t ,
\end{equation}
where $y_{s,t}$ is stratum $s$'s prepaid unpaid principal balance (UPB) in month $t$, in dollars, $n_{s,t}$ its UPB exposure, $x_{s,t}$ the right-hand-side covariates, and $\alpha_{s}$ the stratum fixed effect (295 enter, one absorbed as reference). The term $f(\cdot)$ is a loan-age seasoning spline in mean loan age $a_{s,t}$: a truncated quadratic power basis, a linear term plus $(a - k)_+^2$ at knots $k \in \{12, 24, 36, 60, 84, 120\}$ months, seven basis columns. The remaining terms are $g_{s,t}$, the rate gap in 100 bp units; $b_{s,t}$, the stratum burnout state; and $\phi_t$, the dynamic friction index. The outcome is dollars, not a count. With the log-exposure offset the modeled object is the prepayment rate $y_{s,t}/n_{s,t}$, a nonnegative fraction. So the model is estimated by Poisson pseudo-maximum likelihood, whose score equations remain consistent for the conditional mean without any distributional assumption on $y_{s,t}$. No dispersion statistic is meaningful for this pseudo-likelihood. The 295 stratum fixed effects raise no incidental-parameters problem for the coefficients of interest. Because the model is Poisson, concentrating the fixed effects out of the log-likelihood returns exactly the conditional-likelihood score for the slope vector. The unconditional dummy-variable fit computed here therefore yields the same slope estimates as the conditional estimator, and stays consistent as the number of strata grows with the 36 training months held fixed. The fixed-effect logit or probit does not---its slope estimates are inconsistent at fixed panel length \citep{wooldridge1999}. All that consistency requires is the multiplicative conditional mean $\mathbb{E}[y_{s,t}\mid x_{s,t},\alpha_{s}] = n_{s,t}\exp(\alpha_{s} + x_{s,t}^{\top}\beta)$, with no restriction on the distribution, dispersion, or serial dependence of $y_{s,t}$. All uncertainty statements below come from the two bootstrap schemes of Table~\ref{tab:bootstrap}. The first is a stratum-level cluster bootstrap, robust to arbitrary within-stratum dependence and overdispersion but structurally unable to measure uncertainty in the friction coefficient: $\phi_t$ varies only over time, so resampling strata reproduces the identical friction series in every replication. That burden (and, partially, burnout's) falls to the second, a temporal moving-block bootstrap that resamples the 36 training months in blocks of six with each month's full cross-section intact. Table~\ref{tab:panel} reports the estimation panel's accounting. Both it and Table~\ref{tab:bootstrap} are collected in Appendix~\ref{app:ridge}, because Path A is excluded from every headline range by pre-committed specification and neither table enters a headline quantity. The displayed form \eqref{eq:pathA} is the no-seasonal spec v3 design. The production specification (spec v4) adds eleven calendar-month dummies, and it is the design under which the 121.5\% recovery reported below is computed. But that level is not calendar-attributable (Section~\ref{sec:patha}), and every Path A inferential object stays on spec v3. Table~\ref{tab:hazard-notation} collects the notation shared by the two hazard paths.

The 295 stratum fixed effects absorb each cohort's baseline prepayment speed, so the burnout term captures within-pool adverse selection over time. The burnout regressor is the surviving loan stock, demeaned within stratum and orthogonalized against the loan-age spline so that it is not mechanically collinear with seasoning. This design is the estimation counterpart of \citepos{lesniewski2026} selection identity. The exposure-offset outcome is itself a pool-level, survival-weighted (here UPB-exposure-weighted) hazard. The 295 strata cut the cross-sectional dispersion that drives the identity's negative selection term down to within-cell variation. The demeaned surviving-stock regressor then proxies the residual composition drift the time-invariant fixed effects cannot absorb. An earlier specification used four vintage-year dummies in place of the 295 stratum fixed effects, and returned a burnout coefficient of the wrong sign (+0.76). The full 295-stratum specification, with burnout demeaned within stratum, removes that wrong-signed artifact. The corrected estimate ($-0.13$) carries the theoretically expected sign: a more heavily burned-out pool should prepay more slowly, holding age and rate gap fixed. But the sign statement is a point-estimate statement only. The temporal bootstrap scheme of Table~\ref{tab:bootstrap} centers this coefficient on the opposite sign (mean $+1.23$, median $+0.52$). Under both schemes it is imprecise and not statistically distinguishable from zero. I therefore treat it as a control, not a finding. The fit carries a nominal ridge penalty $\lambda_{\mathrm{ridge}}$ on the full coefficient vector, because 295 fixed effects leave the model ill-conditioned under plain iteratively reweighted least squares. In 20 strata no training prepayment event is recorded at all. The likelihood is then monotone in their fixed effects, and the maximum-likelihood estimates diverge to $-\infty$. I choose $\lambda_{\mathrm{ridge}}$ over the grid $\{10^{-5}, 10^{-4}\}$ on a temporally blocked holdout of all stratum-months from January 2024 onward, so that no within-stratum information crosses the train/holdout split. The penalty is not what stabilizes the fit. At these magnitudes it is a numerical no-op. Both grid points return coefficient vectors identical to machine precision from the IRLS warm start. The production coefficients are therefore the IRLS solution, and the grid selection is inconsequential. The 20 zero-event strata are held at the finite value the warm-started optimizer leaves them at, and the penalty shrinks them only negligibly. The split is temporal over a fixed set of strata, so each stratum's own fixed effect scores its post-January-2024 holdout months. None is estimated on, or transferred to, a stratum absent from the training window. The penalty's operative role is as the fallback for resampled refits in the bootstrap machinery. Nor is the production point an artifact of the device. Dropping the 20 zero-event strata and fitting unpenalized Poisson PML reproduces the spec v3 macro coefficients to within $0.002$ in standardized units (rate gap $+0.6727$ against the spec v3 $+0.6734$; under spec v4 the same check holds to within $0.0022$, Appendix~\ref{app:ridge}). Appendix~\ref{app:ridge} records the verification details under both specifications. It also records a reference-stratum sensitivity of penalized cold-start refits, which motivates the warm-started bootstrap below. The battery has been recomputed under spec v4 and the device is unchanged. Two caveats still attach. The same holdout later serves as the out-of-sample diagnostic, so the RMSE reported below is at least weakly favored by selection. Months never touched by any selection step do not exist in this design, one reason Path A is held to a robustness role.
Estimated on this panel, the production specification --- spec v4, the calendar-month seasonal design adopted at the freeze --- yields a rate-gap coefficient of $+0.65$, a burnout coefficient of $-0.17$, and a friction coefficient of $-0.007$, each with the theoretically expected sign. The rate-gap coefficient is standardized: the training standard deviation of the rate gap is 221 bp, so the equivalent per-100 bp coefficient is $+0.29$. The friction coefficient's shrinkage relative to the prior spec ($-0.037$) is expected under the seasonal design, whose eleven month dummies absorb most within-year variation in the time-only friction index. Point estimates alone overstate what is identified here. Table~\ref{tab:bootstrap} therefore reports cluster-bootstrap uncertainty for all three, including the bias diagnostics. That uncertainty is estimated under the prior spec v3 (see the table note). The rate-gap point estimate sits near the 3rd percentile of its own bootstrap distribution (bootstrap mean $+1.49$, median $+0.76$, estimated bias $+0.82$). That is a signature of resampling-induced upward skew, under which percentile and basic intervals disagree. The percentile interval excludes zero; the bias-reflecting basic interval does not. I therefore rest no claim on interval-based significance for any Path A coefficient. The defensible statements are about direction and imprecision. The rate-gap coefficient's point estimate is positive, and positive in 197 of 198 converged replications (99.5\%). That share is a property of a right-skewed resampling distribution, not a test. The pre-committed bias-respecting test of $H_0\!: \beta_g \le 0$ does not reject under any construction (BCa $p = 0.093$ stratum / $0.412$ temporal; permutation $p = 0.241$; run \texttt{patha\_sign\_test}). The coefficient is therefore directionally consistent with the imported elasticity without being statistically distinguishable from zero. Its magnitude is identified only up to a wide, construction-dependent interval. Burnout and friction are not distinguishable from zero on any construction. All three enter the production simulation as fitted inputs, none as findings, and no result in this paper depends on their point values. Path A's aggregate recovery now carries the interval that pass produces. I resimulate the QT window under each of the 198 converged replications, swapping in the three macro coefficients per draw and holding intercept, age spline, and fixed effects at the production fit. That yields a median of \$922.6 billion and an interquartile range of \$900.1--\$1{,}130.6 billion, but a 95\% percentile interval of $[-\$4{,}459.8, +\$1{,}190.1]$ billion. A joint rerun swaps in each replication's full coefficient vector: intercept, age spline, and at spec v4 the month effects, drawn together with the macros. That rerun returns $[+\$689.0, +\$1{,}217.6]$ billion, far narrower, and it is not adopted. Of its 198 draws, 106 sit at exactly the zero-voluntary-prepayment floor of \$1{,}217.6 billion, among them all 79 rogue-mode draws --- rogue mode meaning a drawn rate-gap coefficient nearer the reference-swap attractor $+1.96$ than the production $+0.67$. That floor is the trapped figure a pool returns when voluntary prepayment is switched off and only scheduled amortization rolls off, measured by a probe through the same coefficient handoff. The committed interval remains the quoted one, because its $-\$4{,}459.8$ billion tail and the joint rerun's floor pile are the same rogue-mode replicates rendered without and with their compensating intercepts. That is one optimizer pathology of penalized refits on resampled panels, not two sampling statements. Measured at spec v4, the same joint rerun fails 1 of 200 replications against spec v3's 2, and its draws avoid the floor entirely (0 of 199). The seasonal design's instability therefore manifests as rogue-mode attraction (73 of 199 draws), not as fit failure. Fifteen replications with extreme burnout or friction draws (the two coefficients the bootstrap classifies as statistically zero) destabilize the 42-month forward compounding entirely, and 83.8\% of draws land above the benchmark. Those two facts pull in opposite directions, and the unfavorable one belongs in the text as much as the favorable one. The same 198 draws have a \emph{mean} of \$598.9 billion, or 78.3\% of the benchmark --- below the rate-inelastic $\beta_1 = 0$ null's 85.7\% and below Path~B's headline 91.3\% --- against a median of 120.6\%. On the share basis the percentile interval is $[-583.2\%, +155.6\%]$ with a resampling standard error of 201.9 points. I report the mean alongside the median for the same reason the coefficient table does. These distributions are severely right-skewed, so no single summary is sufficient; quoting only the median and the fraction above the benchmark would flatter the estimator. Propagated coefficient uncertainty therefore leaves the forward simulation with no useful precision. That is an additional reason the aggregate dollar figure is held to the aggregate-level role assigned below and is not carried in the abstract. Every Path A inferential object remains estimated under the prior spec v3: Table~\ref{tab:bootstrap}, this interval, and the train/holdout decomposition below. That is because bootstrap replications under the spec v4 seasonal design are branch-unstable. The resampled penalized likelihood has a second optimum that attracts replications away from the production branch. On the production panel itself 1 of 50 dispersed starts lands on the production branch, and no start reaches a better penalized objective (run \texttt{pathA\_cold\_starts}). The objective evaluated at the frozen coefficient vector is the best value any of 150 starts attains, and every non-production terminus sits at least seven orders of magnitude above it. That is an optimizer stall of a strictly convex objective. Until the resampled instability is stabilized, the production specification carries no valid uncertainty of its own. Path A is excluded from the paper's headline figures for this reason, retained as an aggregate-context, non-corroborating exhibit. Simulating forward, this model recovers \$928.9 billion, or 121.5\%, of the \$764.7 billion empirical benchmark under the production calendar-month specification (spec v4). The forward simulation drains each cohort's balance deterministically each month by balance times hazard, with fractional prepayment settling directly to simulated SOMA roll-off and no settlement-lag convolution. That specification was adopted at the pre-submission freeze under pre-committed parity gates that reproduce the committed seasonal robustness artifact exactly. The prior specification is preserved as a spec-v3 exhibit.

Three diagnostics temper this aggregate result, because they bear directly on how much weight the \$928.9 billion figure can carry. First, the holdout fit, evaluated under spec v3, is weak in both weightings. Re-scoring the frozen spec v4 fit on the same January-2024-onward holdout, without refitting, gives 37.9 unweighted and 3.0 exposure-weighted percentage points against spec v3's 37.8 and 2.5, so the seasonal terms do not improve the out-of-sample fit. The units matter here. The unweighted RMSE of approximately 38 percentage points of annualized stratum-month CPR is dominated by small-exposure strata with unstable realized rates, while the exposure-weighted RMSE, the weighting the fit itself uses, is 2.5 percentage points against an empirical mean CPR of 5.14\% (ABM basis), half the mean. Neither weighting supports out-of-sample predictive power at the stratum-month level, and the ridge grid's two candidate fits are numerically identical on this holdout, so no selection hinges on the weighting. A flexible-learner ceiling test bounds how much of this weakness is functional form. A gradient-boosted Poisson learner given the same covariate information on the identical pre-committed holdout improves the exposure-weighted RMSE materially, to 0.9 points against the frozen spec v4's 3.0. It leaves the unweighted RMSE essentially unmoved (37.6 against 37.9) and out-of-sample $R^2$ at zero ($+0.003$). A penalized Poisson GLM on the exact v4 design does no better than the production fit (2.6 and 38.1 points). Read against the run's pre-committed thresholds, flexibility buys recalibration of the high-exposure cells but no cross-cell predictive power in the 2024-onward regime. The landmarking-with-recalibration strand receives the same test. An exposure-weighted isotonic map fit on the 2023 inner-validation year only, applied to the frozen spec v4 holdout predictions, improves the exposure-weighted RMSE to 2.53 points against the frozen 3.04, short of the pre-committed 0.80$\times$ threshold. It leaves the unweighted RMSE at 37.8 and out-of-sample $R^2$ at zero ($-0.007$). The pre-committed verdict is no material change: a monotone level recalibration fit in the training regime does not transfer predictive power across the 2024 boundary. Second, the recovery is concentrated where the model is unsupervised. The decomposition is measured under spec v3 and was not recomputed at adoption. Of the spec v3 fit's \$915.1 billion, \$331.0 billion accrues in the 19 QT months overlapping the training window (June 2022--December 2023), against \$337.5 billion empirical (98\%). A further \$584.1 billion accrues in the 23 post-training simulation months from January 2024 onward, against \$427.3 billion empirical (137\%); 21 of those months have performance data, and the empirical comparison uses those. Components are independently rounded, and the overshoot lives almost entirely out of sample. Third, the raw CPR correlation against the empirical path is $r = -0.444$ at zero lag, more negative even than the ABM's $-0.318$, and \eqref{eq:pathA} omits calendar-month seasonality, one of the four factors of \citepos{richard1989} canonical prepayment function and a natural first suspect for a monthly path fit this poor. Re-estimating with eleven calendar-month dummies confirms strong, correctly signed seasonality (spring and early-fall log-effects up to $+0.43$ against a January base) and improves the simulated path's contemporaneous correlation from $-0.444$ to $-0.378$, with the peak moving from lag $-3$ to $-2$. That timing improvement is calendar-attributable, and I test it. Permuting the assignment of the twelve estimated month log-effects to calendar months preserves their multiset exactly, and therefore every level channel they drive, while destroying only their alignment with calendar time. I score the permutation on the conjunctive criterion the claim actually states: a lag-0 correlation better than $-0.378$ \emph{and} a peak at lag $-2$. Scrambled assignments win 1 of 64 draws, a one-sided $p$ of 0.031. An independent sixteen-draw family at a different seed also returns a single win (\texttt{pathA\_seasonal\_permutation}). The two components scored separately, 11 and 8 of 64, are diagnostics only: each necessarily exceeds the joint count, so reading either alone understates a conjunctive claim. The recovery figure is a different matter, and I no longer credit it to the calendar terms. Across those same 64 scrambles the recovery spans 121.14--122.29\%, a range of 1.15 points that contains the adopted 121.5\%, so the level is calendar-invariant: permutation holds the multiset of month effects fixed, and it is the multiset, not its alignment, that sets the level. Nor is the spec-v4 total decomposable as production plus seasonality in the first place. The seasonal refit with all twelve month multipliers zeroed moves the recovery by $+\$80.27$ billion on its own, and restoring the multipliers claws back $-\$66.44$ billion of that, against a net reported seasonal effect of $+\$13.83$ billion. The timing improvement is real but far from sufficient: the named first suspect accounts for only a modest share of the negative path fit. A model that recovers 121.5\% of the dollar benchmark while mismatching the monthly path this badly may be recovering the aggregate partly through offsetting errors, over-predicting prepayment in some cohort-months and under-predicting in others in a way that nets out near the benchmark without the underlying mechanism being correctly timed. I therefore treat Path A's recovery as evidence about the aggregate level only. No timing claim rests on Path B either: its apparent levels advantage does not survive detrending (Section~\ref{sec:pathb}), and every claim this paper makes for Path B is a level-accuracy claim. Table~\ref{tab:pathadiag} collects Path A's diagnostics in one place.

\begin{table}[H]
\centering
\caption{Path A diagnostic summary. An aggregate-context, non-corroborating exhibit: Path A is excluded from the paper's headline figures by pre-committed specification, and no cell of this table enters a headline quantity. Values repeat committed figures derived in this section and Appendix~\ref{app:theil}.}
\label{tab:pathadiag}
\footnotesize
\begin{threeparttable}
\begin{tabular}{@{}p{4.6cm}p{9.6cm}@{}}
\toprule
Diagnostic & Value \\
\midrule
Aggregate recovery \emph{(spec v4)} & \$928.9 billion $=$ 121.5\% standalone; 112.4\% shared basis. Order of magnitude coincides with Path B; the sign does not corroborate it, the coefficient being statistically indistinguishable from zero (row below and Section~\ref{sec:patha}) \\
Fit diagnostics \emph{(spec v4)} & exposure-weighted holdout RMSE 3.0 points; Theil $U_1$ 0.431, $U_2$ 1.010; terminal cumulative error $+\$164.1$ billion \\
\bottomrule
\end{tabular}
\begin{tablenotes}
\footnotesize
\item \emph{Notes:} Every row above is measured under the production calendar-month specification (spec v4), so this table is single-specification and no cell need be read against a neighbour measured under a different design. The spec v3 objects the table previously carried in its body are recorded here instead, demoted because they were not recomputed at adoption: the \emph{refit-and-resimulate distribution} --- point \$915.1 billion; median \$922.6 billion; \emph{mean \$598.9 billion $=$ 78.3\% of benchmark}; IQR \$900.1--\$1{,}130.6 billion; 95\% percentile interval $[-\$4{,}459.8, +\$1{,}190.1]$ billion ($[-583.2\%, +155.6\%]$ on the share basis); 83.8\% of draws above benchmark --- the \emph{train vs post-training split}, \$331.0 billion (98\% of empirical) vs \$584.1 billion (137\%); and the spec v3 counterpart of the holdout RMSE, 2.5 points against the spec v4 row's 3.0. The resampling distribution has since been rerun jointly at both specifications (Section~\ref{sec:patha}), and the committed spec v3 interval remains the quoted one --- the joint rerun is narrower for the mechanical reason stated there. The spec v3 point (\$915.1 billion) is recorded here, rather than beside the spec v4 aggregate, so that the median is not read as centred on a level measured under a different design. The 95\% interval spanning zero is the stated reason for the pre-committed headline exclusion; the flexible-learner and isotonic comparators of this section bound what refitting the same information could recover.
\end{tablenotes}
\end{threeparttable}
\end{table}

\begin{figure}[H]
\centering
\includegraphics[width=0.9\textwidth]{fig2_ccf}
\caption{Cross-correlation functions between each estimator's simulated CPR path and the empirical SOMA-derived path, lags $-6$ to $+6$. Under the implemented convention (Section~\ref{sec:pathb}), a peak at negative lag means the empirical series leads and the simulated path trails. Path B's $-3$/$-4$ peak and Path A's and the ABM's negative contemporaneous correlations summarize the same fact from Table~\ref{tab:estimators}: no estimator shows positive contemporaneous alignment, and once the window's shared trend is removed no estimator's co-movement is distinguishable from zero under the moving-block bands.}
\label{fig:ccf}
\end{figure}

A moving-block bootstrap (block length 6, 4,000 replications) puts the 95\% interval on the peak-lag correlation at $[-0.21, +0.49]$ and on the lag-0 correlation at $[-0.13, +0.44]$. With seven lags searched, individual band-point differences in peak $r$ carry no evidential weight. These intervals are not sensitive to the block-length choice in any verdict-relevant way. At block lengths 4 and 8 the peak-lag endpoints move to $[-0.27, +0.39]$ and $[-0.22, +0.51]$, and the lag-0 endpoints to $[-0.11, +0.45]$ and $[-0.17, +0.42]$. Every interval still contains zero.

A Theil decomposition formalizes where each estimator's monthly error lives. It is a fit diagnostic, not a timing credential. On levels, the ABM's mean-squared error is 68\% bias share (the level failure already visible in its aggregate). Path B's is 11.3\% bias and 82.5\% variance share on the corrected amortization basis (run \texttt{coupon\_convention\_amortization}). Its residual monthly error is under-dispersion against a seasonal empirical path, not level bias. Path A is intermediate (26.5\% bias, 63.0\% covariance share). After linear detrending the variance share still dominates for Path B (93.7\%). On first differences, Theil's $U_2$ is 1.23 for the ABM, 1.01 for Path A, and 1.00 for Path B. No estimator beats a naive no-change forecast of month-over-month CPR changes. That restates the timing concession above in forecast-accuracy units. Table~\ref{tab:theil} collects the full decomposition by transform. Its cumulative dollar-error profiles have terminal values that reproduce Table~\ref{tab:estimators}'s aggregates by construction.

A direct input-timing diagnosis rules out the dating of the rate input. It narrows the mechanism, but does not isolate it. In monthly \emph{changes}, the simulated CPR responds to the contemporaneous rate change with the correct sign and near-mechanical strength ($\Delta$-correlation $-0.94$ at lag 0). The empirical CPR's month-to-month changes show no detectable rate response ($+0.25$, inside the $\pm 0.31$ zero band at $n = 41$). At monthly frequency the empirical path's variation is dominated by non-rate factors the simulation does not model, calendar seasonality above all. One construction asymmetry belongs in this reading as well. The empirical CPR back-out retains the income-scaled curtailment component (Section~\ref{sec:method-benchmark}); the simulated CPR paths exclude it, since curtailment enters only as a common dollar flow in the accounting layer. The two series being correlated therefore differ by a small, slowly-varying component as well as by seasonality.

\begin{table}[H]
\centering
\caption{Danish-leg sensitivity to the buyback discount rate. The diagnostic executed on the Danish-level Berger-recalibrated leg (mean Danish CPR 3.39\%), which is the bracketing anchor of Section~\ref{sec:abm-danish} rather than the production U.S.-intercept leg. Neither leg consumes the present-value classifications the discount rate enters, so the zero deltas hold on the standalone, shared, and printed bases alike; under the production anchor that independence is structural rather than run-verified. The baseline CPR path reproduces the frozen parquet for the leg that ran exactly. The zeros bound the hazard channel only: the leg's dollar conversion credits retired balance at par, an un-priced channel sized in Section~\ref{sec:pathb}.}
\label{tab:discount}
\begin{tabular}{lccc}
\toprule
Discount-rate shift & PV classifications flipped & $\Delta$ CPR path (pp) & $\Delta$ trapped (\$B) \\
\midrule
0 bp (production) & -- & 0.00 & 0.00 \\
$-25$ bp & 2{,}470 of 1{,}683{,}082 & 0.00 & 0.00 \\
$-50$ bp & 7{,}805 of 1{,}683{,}082 & 0.00 & 0.00 \\
$-75$ bp & 17{,}493 of 1{,}683{,}082 & 0.00 & 0.00 \\
$-100$ bp & 32{,}488 of 1{,}683{,}082 & 0.00 & 0.00 \\
$-125$ bp & 55{,}210 of 1{,}683{,}082 & 0.00 & 0.00 \\
$-150$ bp & 92{,}495 of 1{,}683{,}082 & 0.00 & 0.00 \\
\bottomrule
\end{tabular}
\end{table}

The decomposition also carries its time axis. The monthly marginal path is the central-minus-null difference of the two committed simulated legs. Its sum reproduces the $+\$70.3$ billion bit-exactly. That path is positive in all 42 window months. It splits 21.3\%/41.9\%/36.8\% across the window's thirds and peaks at \$2.29 billion in November 2023, the window's eighteenth month. By coupon bucket the terminal marginal divides 38.0\%/43.7\%/18.3\% across the $<3.0\%$/3.0--4.0\%/$\geq 4.0\%$ groups, on 40.3\%/40.4\%/19.3\% of sample balance. Figure~\ref{fig:marginaltiming} draws the stack.

\begin{figure}[H]
\centering
\includegraphics[width=0.94\textwidth]{fig12_marginal_timing.png}
\caption{Monthly lock-in marginal by coupon bucket. Stacked areas are the per-bucket monthly marginals from the month-by-cell capture of the group-ablation engine; the dark line is the aggregate central-minus-null path, whose sum reproduces the $+\$70.3$ billion in-sample marginal bit-exactly. Dotted rules mark the window's thirds; the marker is the November 2023 peak. The stack differs from the aggregate by the decomposition's 1.03\% terminal additivity residual, asserted in the generator.}
\label{fig:marginaltiming}
\end{figure}

\section{Path A Estimation Panel, Coefficient Uncertainty, and Ridge Device}\label{app:ridge}

The replication package verifies three facts about the ridge device of Section~\ref{sec:patha}, under both specifications: the spec v3 records stand, each fact has been recomputed under spec v4, and the device is unchanged. First, at production magnitudes the penalty is a numerical no-op on the full panel. The two grid points ($\lambda_{\mathrm{ridge}} \in \{10^{-5}, 10^{-4}\}$) return coefficient vectors identical to machine precision from the IRLS warm start (maximum parameter difference $4.5\times10^{-17}$ in the replication artifact; $8.4\times10^{-17}$ under spec v4). Grid selection is therefore inconsequential, and the production coefficients are the IRLS solution to machine precision; ``mildly shrunk'' would overstate an effect that is nil. Second, ridge on the full coefficient vector with one stratum absorbed as reference is not invariant to the reference choice. In cold-start refits on resampled panels, swapping the reference stratum can move the standardized rate-gap coefficient by more than its own magnitude --- an optimizer-path artifact of penalized refits with the 20 inestimable fixed effects. The bootstrap machinery accordingly warm-starts every replication and discards non-converged draws. The warm-started spec v4 swap refit is reference-invariant to $7.3\times10^{-14}$, consistent with the warm-start policy that concession motivated. That invariance is not evidence that the cold-start sensitivity is gone. Third, the production point itself is not an artifact of these choices. A refit that drops the 20 zero-event strata (423 of 10,176 cells) and fits unpenalized Poisson PML converges, is exactly invariant to the reference choice, and reproduces the spec v3 macro coefficients to within $0.002$ in standardized units: rate gap is $+0.6727$ against the spec v3 $+0.6734$, and friction $-0.0373$ against $-0.0371$, third-decimal agreement. Burnout is $-0.1317$ against $-0.1301$; the larger burnout gap reflects that dropping the 423 cells changes the burnout regressor's within-stratum demeaning and age-orthogonalization sample. At spec v4 the identical 20 strata and 423 cells drop out, since the month dummies do not change the zero-event set. The same refit is again exactly reference-invariant ($1.1\times10^{-14}$) and reproduces the spec v4 macro coefficients to within $0.0022$: rate gap $+0.6463$ against $+0.6469$, friction $-0.00715$ against $-0.00711$, burnout $-0.1734$ against $-0.1713$, against a tolerance pre-committed at $0.005$, loosened from v3's $0.002$ because the month dummies absorb time variation the friction coefficient formerly carried.

Table~\ref{tab:panel} reports the estimation panel this fit is built on. Table~\ref{tab:bootstrap} reports the coefficient uncertainty the two bootstrap schemes of Section~\ref{sec:patha} return. Both sit here because Path A is an aggregate-context, non-corroborating exhibit that pre-committed specification excludes from every headline range. No cell of either table enters a headline quantity.

\begin{table}[H]
\centering
\small
\begin{threeparttable}
\caption{Path A estimation panel.}
\label{tab:panel}
\begin{tabular}{@{}lrr@{}}
\toprule
 & Training & Holdout \\
\midrule
Performance months & Jan.\ 2021--Dec.\ 2023 & Jan.\ 2024--Sep.\ 2025 \\
Stratum-months & 10,176 & 6,077 \\
Prepaid UPB (PML outcome) & \$650.0B & \$131.3B \\
UPB exposure (month-sums) & \$72.5T & \$37.4T \\
\bottomrule
\end{tabular}
\begin{tablenotes}[flushleft]\footnotesize
\item The panel aggregates the full Freddie Mac 2017--2021 origination universe (8.84 million active loans at the window's peak month, not the 75,000-loan sample that Section~\ref{sec:pathb} simulates) to 296 strata, and carries no sampling weights: exposure and the outcome are raw dollars, so the pseudo-maximum-likelihood fit is UPB-weighted by construction (a dollar of prepayment counts the same in any stratum-month). The temporal split at January 2024 is the ridge-selection holdout of the text; the performance window is January 2021 through September 2025.
\end{tablenotes}
\end{threeparttable}
\end{table}

\begin{table}[H]
\centering
\footnotesize
\setlength{\tabcolsep}{4pt}
\begin{threeparttable}
\caption{Path A coefficient uncertainty under two bootstrap schemes (estimated under the prior specification, spec v3; see note).}
\label{tab:bootstrap}
\begin{tabular}{@{}lrrrrccr@{}}
\toprule
Coefficient & Point & Mean & Median & SE & Percentile CI & Basic CI & Draws $\le 0$ \\
\midrule
\multicolumn{8}{@{}l}{\textit{Stratum-cluster scheme (198 of 200 replications converged)}} \\
Rate gap (standardized) & $+0.673$ & $+1.491$ & $+0.763$ & $1.152$ & $[+0.60,\,+4.11]$ & $[-2.76,\,+0.74]$ & 0.5\% \\
Burnout (orthogonalized) & $-0.130$ & $-0.056$ & $-0.130$ & $1.363$ & $[-1.76,\,+1.43]$ & $[-1.69,\,+1.50]$ & 75.3\% \\
Friction & $-0.037$ & $+0.352$ & $+0.015$ & $0.506$ & $[-0.25,\,+1.38]$ & $[-1.46,\,+0.18]$ & 48.0\% \\
\midrule
\multicolumn{8}{@{}l}{\textit{Temporal moving-block scheme, block length 6 months (196 of 200 converged)}} \\
Rate gap (standardized) & $+0.673$ & $+3.313$ & $+2.556$ & $3.845$ & $[+0.44,\,+12.18]$ & $[-10.84,\,+0.91]$ & 0.5\% \\
Burnout (orthogonalized) & $-0.130$ & $+1.229$ & $+0.523$ & $2.586$ & $[-3.16,\,+6.16]$ & $[-6.42,\,+2.90]$ & 40.8\% \\
Friction & $-0.037$ & $+1.050$ & $+0.567$ & $2.802$ & $[-3.71,\,+7.69]$ & $[-7.76,\,+3.63]$ & 37.2\% \\
\bottomrule
\end{tabular}
\begin{tablenotes}[flushleft]\footnotesize
\item Stratum-cluster scheme: the observed vintage--coupon--FICO--LTV strata are resampled with replacement, each resampled stratum receives a fresh fixed-effect label so duplicated strata contribute independent fixed effects, and per-replication standardization (including the burnout orthogonalization, which is refit inside every replication) is rescaled to production standardized units. This scheme is structurally silent about the friction coefficient: the friction index varies only over time, so every replication sees the identical series, and the stratum-scheme friction row understates that coefficient's uncertainty by construction. Temporal scheme: the 36 training months are resampled in moving blocks of six with each month's cross-section kept intact, same refit and rescaling; it is the informative scheme for friction and, partially, burnout; its friction SE (2.80) is roughly $5.5\times$ the stratum scheme's (0.51). The ridge penalty is held at the production $\lambda_{\mathrm{ridge}} = 10^{-4}$ rather than re-selected per replication (re-selection would bootstrap a different, coefficient-plus-model-selection estimator). Replications returning non-finite or implausible ($|\beta| > 20$) coefficients are counted as failures, not draws; the point refit reproduces the production coefficients to four decimals on the same 10,176 training cells. Coefficients are in production standardized units (training SD of the rate gap: 221 bp). The stratum-scheme burnout median coincides with the point estimate at reported precision only (unrounded: median $-0.1297$, point $-0.1301$). Mean, median, and the implied bias (mean $-$ point) are reported because both schemes' rate-gap distributions are heavily right-skewed; under that skew the percentile and basic 95\% intervals disagree about zero, so interval-based significance is construction-dependent and the text relies on sign stability only; but the bootstrap sign-share is itself percentile logic, and the pre-committed bias-respecting test of $H_0\!: \beta_g \le 0$ (run \texttt{patha\_sign\_test}; BCa with the 296-refit jackknife acceleration this note previously deferred, plus a month-permutation leg) does not reject under any construction ($p = 0.093$/$0.412$/$0.241$), so the 99.5\% positive share is a description of a right-skewed resampling distribution, not a sign test, and no sign claim for this coefficient survives.
\end{tablenotes}
\end{threeparttable}
\end{table}

\section{Dynamic-Fit Decomposition}\label{app:theil}

Table~\ref{tab:theil} collects the Theil fit diagnostics summarized in Section~\ref{sec:pathb}: $U$ statistics, the bias/variance/covariance decomposition of each estimator's mean-squared error on three transforms of the monthly CPR paths, and the terminal cumulative dollar errors. The $\beta_1 = 0$ null's cumulative-error profile follows from the committed legs by arithmetic. Its terminal cumulative error against the empirical path is $-\$16.6$ billion, versus Path B's $+\$53.8$ billion. The null under-runs the empirical roll-off by approximately the lock-in marginal, by construction. The full monthly profile is committed with the marginal-timing data of Section~\ref{sec:identification}.

\begin{table}[H]
\centering
\caption{Theil dynamic-fit decomposition for the three estimators' monthly CPR paths over the 42-month QT window. Fit diagnostics, not timing credentials: no estimator's co-movement survives detrending (Section~\ref{sec:pathb}), and $U_2 > 1$ means worse than a naive no-change forecast of month-over-month changes.}
\label{tab:theil}
\footnotesize
\begin{threeparttable}
\begin{tabular}{@{}lrrr@{}}
\toprule
 & ABM & Path A & Path B \\
\midrule
Theil $U_1$, levels & 0.412 & 0.438 & 0.271 \\
Theil $U_1$, detrended & 0.807 & 0.839 & 0.980 \\
Theil $U_2$, first differences & 1.226 & 1.009 & 1.003 \\
\midrule
\multicolumn{4}{@{}l}{\textit{MSE shares, levels (\%)}} \\
\quad Bias & 68.1 & 30.6 & 11.3 \\
\quad Variance & 0.4 & 10.5 & 82.5 \\
\quad Covariance & 31.5 & 58.8 & 6.3 \\
\multicolumn{4}{@{}l}{\textit{MSE shares, first differences (\%)}} \\
\quad Bias & 0.1 & 0.3 & 0.1 \\
\quad Variance & 14.0 & 71.4 & 96.8 \\
\quad Covariance & 85.8 & 28.3 & 3.1 \\
\multicolumn{4}{@{}l}{\textit{MSE shares, linearly detrended (bias $\equiv 0$, \%)}} \\
\quad Variance & 0.5 & 43.6 & 93.8 \\
\quad Covariance & 99.5 & 56.4 & 6.2 \\
\midrule
Terminal cumulative error (\$B) & $-673.8$ & $+164.1$ & $+53.8$ \\
\bottomrule
\end{tabular}
\begin{tablenotes}[flushleft]\footnotesize
\item $U_1 \in [0,1]$, where zero is a perfect fit; $U_2$ is relative to the naive no-change forecast. Path A and Path B columns are computed against the hazard-side empirical CPR back-out at the pool's note-rate WAC (the corrected amortization basis; run \texttt{coupon\_convention\_amortization}, committed 2.5\%-convention values in \texttt{theil\_data.json}); the levels-bias shares are the cells the correction moves ($+4.6$pp for Path B), the timing rows move at or below a thousandth. The ABM column retains its committed basis --- the converted ABM reconstruction failed its own replay gate and is not quoted (the ABM's cumulative-error leg uses its own committed back-out). Terminal cumulative errors equal each estimator's trapped-liquidity aggregate minus the empirical \$764.7 billion on its own accounting basis, reproducing Table~\ref{tab:estimators} by construction. Artifact \texttt{theil\_data.json}; inputs and correlations are gated to the committed run manifests.
\end{tablenotes}
\end{threeparttable}
\end{table}

Grouped calibration localizes the same verdict (run \texttt{grouped\_calibration}). I bin the frozen Path A cell predictions into five exposure-weighted prediction-CPR quintiles, with edges fit on the training window only. I then cross them with calendar half-years. The training halves calibrate at a pooled predicted-to-realized ratio of 1.000, exact by construction of the exposure-offset fit. Per-half ratios span 0.80--1.17. The post-2024 holdout instead drifts into monotonically worsening under-prediction: 0.64, 0.57, 0.44, 0.43 across 2024H1--2025H2, pooled predicted 2.24 against realized 4.21 points. The departure is temporal drift, not a reliability-slope failure. Pooled across periods the four upper prediction bins calibrate at ratios 0.94--1.11, and the shortfall concentrates in the lowest-prediction bin (43.5\% of exposure), where the 2024--25 cells accumulate. That is the finer-resolution form of the $U_2 > 1$ diagnosis above. This is a cell-panel CPR object, distinct from the cumulative dollar-runoff split of Section~\ref{sec:patha}. Like every Path A diagnostic, it moves no headline figure.

\section{Composition Shift: Estimation Sample versus SOMA Book}\label{app:composition}

Table~\ref{tab:composition} reports a formal covariate-shift analysis between the Freddie Mac data behind the hazard paths and the SOMA book the benchmark spans, on the dimensions both sides carry. The book side is the New York Fed's CUSIP-level SOMA holdings. The same production code that builds the ABM cohorts parses them, anchor-gated to the committed marginals (book WAC 2.49\%, Ginnie share 20.4\%, term split, vintage-group shares). The API also served the historical June 29, 2022 as-of, whose composition is consistent with those anchors. The artifact reports that as-of as a non-gated secondary block. The shifts on coupon, vintage, and agency are large --- the expected result, disclosed in prose throughout the paper. Each shift carries a committed, signed bound as its consequence, listed in the table's final column. The joint agency$\times$vintage cells settle the one netting question the marginals leave open. Of the 2022 vintage's 23.1\% of book face, 20.7\% is conventional and 2.4\% is the Ginnie intersection. Vintage and agency marginals therefore cannot be added without double-counting that cell. Geography, loan size, FICO, and LTV do not exist at the CUSIP level of the Fed's published holdings, which carry par value and security description only. Panel (b) supplies the deliverable substitute: the Freddie and Fannie loan-level samples agree to negligible shift on exactly those dimensions. The hazard specification's exposure to the loan-level covariates is bracketed directly by the printed covariate-block endpoints (99.2\%/92.2\% shared-basis, Section~\ref{sec:pathb}).

\begin{table}[H]
\centering
\caption{Composition comparison. Panel (a): the two Freddie objects the hazard paths use --- the estimation universe Path A is fitted on (exposure-weighted at the June 2022 window open) and, indented beneath it, the 75,000-loan draw Path B simulates (by count and on its surviving balance at the window open) --- against the SOMA agency-MBS book (CUSIP-level holdings, anchor-gated). The three are distinct populations and none is mixed within a row; only the unindented rows carry a PSI against the book. Panel (b): cross-agency stability of the loan-level samples on the dimensions the book cannot supply. PSI is the population stability index $\sum_i (p_i - q_i)\ln(p_i/q_i)$; values above 0.25 are conventionally read as large shift, below 0.10 as negligible.}
\label{tab:composition}
\footnotesize
\begin{threeparttable}
\begin{tabular}{@{}p{2.1cm}p{4.7cm}p{4.7cm}p{3.6cm}@{}}
\toprule
Dimension & Freddie side (\%) & Book (\%) & PSI; committed consequence \\
\midrule
\multicolumn{4}{@{}l}{\emph{(a) Freddie side vs SOMA book: estimation-universe rows, with Path B's draw indented beneath}} \\
Agency & UMBS 100.0; Ginnie 0.0 & UMBS 79.6; Ginnie 20.4 & 2.54; static bound \$20--47B, overlay differential $+\$38.5$B \\
Vintage & pre-2017: 0.0; 2017--19: 12.5; 2020: 37.3; 2021: 50.2; 2022: 0.0 & 10.6; 6.0; 16.3; 43.9; 23.1 & 4.33; vintage residual bound \$11.7 billion \\
\quad Path B draw & \emph{count}: 0.0; 60.0; 20.0; 20.0; 0.0 \newline \emph{surviving balance}: 0.0; 27.0; 30.6; 42.3; 0.0 & & no committed draw-vs-book vintage PSI; equal allocation per origination quarter, never reweighted (Appendix~\ref{app:params}) \\
Coupon & $<3.0\%$: 18.0; 3.0--4.0\%: 68.2; $\geq 4.0\%$: 13.8 & 73.9; 18.7; 7.3 & 2.29; full-book reweighting 107.0\% $\to$ 109.1\% \\
\quad Path B draw & \emph{count}: 7.1; 34.7; 58.2 \newline \emph{surviving balance}: 14.0; 58.8; 27.2 & & 3.38 on the draw's origination-balance shares; draw mean note rate 3.86\%; the reweight moves the level, not the marginal \\
Term & 30-year collateral (15-year face excluded from hazard cohorts, disclosed in Section~\ref{sec:limitations}) & 30yr 90.7; 15yr 9.1; other 0.2 & book-only dimension \\
\midrule
\multicolumn{4}{@{}l}{\emph{(b) Freddie vs Fannie loan-level samples (75{,}000 loans each, origination-balance weighted)}} \\
State & \multicolumn{2}{l}{PSI 0.018} & negligible \\
Loan size & \multicolumn{2}{l}{PSI 0.000} & negligible \\
FICO & \multicolumn{2}{l}{PSI 0.005} & negligible \\
LTV & \multicolumn{2}{l}{PSI 0.000} & negligible \\
\bottomrule
\end{tabular}

\end{threeparttable}
\end{table}

\noindent{\footnotesize \emph{Notes to Table~\ref{tab:composition}.} Three populations appear in panel (a), and no row mixes them. The unindented rows are the hazard estimation universe, exposure-weighted at the window open; the indented rows are the 75,000-loan Path B draw, by count and on the surviving balance the aggregation applies; the Book column is SOMA face. All three sides use rounded half-point coupon buckets. But the two Freddie sides bucket borrower note rates while the Book column buckets security pass-through coupons --- a wedge of $+0.80$ points --- so the buckets are aligned in width, not in object. Within the draw the $\geq 4.0\%$ share moves with both conventions, and the two must be read separately. On count weights, rounding to half points raises it from 47.0\% to 58.2\%. Holding the rounding fixed, the weighting lowers it from 58.2\% by count to 53.7\% on origination balance and 27.2\% on surviving balance. The WAC pair is mixed-basis, as Section~\ref{sec:pathb} records: the draw's 3.86\% is the unweighted mean of borrower note rates, the book's 2.49\% the face-weighted mean of security pass-through coupons. Every shift is signed in the same conservative direction. The out-of-sample segments prepaid faster, so trapped liquidity is if anything overstated (Section~\ref{sec:pathb}). That signing covers observed speeds in the segments outside the draw's support. It does not cover what the draw's composition within 2017--2021 does to the floor's bind share, which no committed run scores (Appendix~\ref{app:params}). Panel (b) uses the committed Fannie replication sample under its aggregates-only license posture. Artifact \texttt{composition\_shift\_results.json}.\par}

\nocite{cfpb_atr,frbny_soma,fred_actliscou,fred_wshomcb,fred_dspic96,fred_umcsent}

\section{Floor-Sweep and Contamination Mechanics}\label{app:floormech}

This appendix collects the mechanics behind the calibration box of Section~\ref{sec:robustness-floor}. It covers the floor-level sweep, the marginal as a curve in the imported elasticity, the separation of the two refinance contaminations behind the off-window floor anchors, the out-of-sample anchor construction, the calibration catalogue for the headline marginal row, and the inference ladder on the floor read (Table~\ref{tab:ladder}). The three exhibits themselves are read where the headline is set, in Section~\ref{sec:hazard} (Tables~\ref{tab:floorband}, \ref{tab:lowband}, and~\ref{tab:oosfloor}).

Table~\ref{tab:floorband} reports the lock-in marginal at every cell of the resulting calibration box. Three facts shape the reading. First, the pre-committed stability threshold ($\pm 2$ points across the empirically plausible 3--5\% floor range) fails: at the central elasticity the marginal spans $+10.8$ to $+5.5$ points across that range. The mechanism is structural. The hazard specification takes a maximum against the floor, so wherever the floor binds the elasticity is inert. Raising the floor then mechanically crowds out the lock-in channel. Second, both pre-committed monotonicities hold: the marginal decreases in the floor within every band row and increases in the band at every floor. The box maximum lands at the predicted corner, $+13.2$ points ($+\$100.7$ billion) at a 2\% floor on the 7.7\% band edge. Third, every cell is positive, the minimum $+2.1$ points at a 6\% floor on the 5.5\% band edge. Every marginal in the box is also basis-invariant, because the shared-accounting netting is common to the central and null legs of each cell (Section~\ref{sec:robustness-hybrid}). That positivity is a consistency check on the specification, not an identifying result, since the window's rate configuration fixes the sign before any floor or band is chosen (Section~\ref{sec:identification}). The verification reading of Section~\ref{sec:pathb} accordingly weakens from a point to a bounded interval. The lock-in elasticity's identified contribution is $+9.2$ points at the production calibration, $+5.0$ to $+12.6$ points across the empirically plausible floor range, and $+2.1$ to $+13.2$ points across the full box. The off-window re-measurement below places the headline reading at $+4.3$ to $+6.8$ points. The aggregate level remains floor-dominated throughout.

The box sweeps the elasticity only over its published band, and that band's lower edge is a literature boundary, not an evidential one. The in-sample estimates cannot reject zero for the elasticity (Section~\ref{sec:pathb}). A reader who weighs that non-rejection over the imported band therefore needs the marginal at values the box never visits. A pre-committed extension (run \texttt{band\_low\_extension}) re-runs the paired legs at quarterly-decline values $\delta \in \{0, 1.0, 2.0, 3.0, 3.25, 4.0, 5.0\}\%$ below the band, at both anchor floors. Its grid, its parity gates against the committed production and off-window cells, and its monotonicity check were fixed in the script header before any run. Table~\ref{tab:lowband} reports the resulting curve beside the band's own cells. The curve is monotone at both floors and passes through zero at $\delta = 0$ by construction. At the off-window floor it is also concave, through the censoring mechanism of Section~\ref{sec:identification}: raising the response crowds the hazard into the floor (the bind share climbs from 35.8\% to 73.5\% along that curve), so each increment of elasticity buys less marginal. At half the adopted central, $\delta = 3.25\%$, the marginal is $+5.1$ points ($+\$39.1$ billion) at the in-sample floor and $+3.5$ points ($+\$26.5$ billion) at the headline off-window floor. This curve prices the dependence on the imported elasticity rather than identifying it. How small the marginal could be, if the moving response is weaker than the literature band, is read off the curve; the identified content remains the bounded range stated in Section~\ref{sec:identification}.


Those two checks vary the floor's level and its functional form. A third pre-committed check interrogates the level's provenance: whether the 4\% figure is an artifact of being measured in-window on the locked-in cohorts the marginal later decomposes. Two runs carry it, \texttt{out\_of\_window\_floor} and \texttt{oos\_identification}, and spec, anchor definitions, acceptance thresholds, and parity gate were fixed in each script header before its runs, the ex-ante convention the replication package's ledger records. Run in-window on the deeply out-of-the-money 2023--24 discount cohorts, the identical exposure-weighted-CPR procedure reproduces the production floor at 3.8--3.9\% annual CPR, flat across seasoning cuts and validating the measurement where it is used. Applied out-of-window to pooled 2017--2019 discount-cohort turnover, before both the QT window and the 2020--21 refinancing wave, it recovers a materially higher 6.1\%. That pooled reading, however, does not survive its own diagnostics, and I do not take it as the operative off-window anchor.

Two contaminations are separable, and separating them is what changes the reading. The first is seasoning. A floor meant to capture strictly involuntary turnover---death, divorce, relocation---should be approximately flat in loan age, but the pooled 2017--2019 out-of-the-money turnover \emph{rises} monotonically with it: 3.68\% below twelve months, 5.78\% at twelve to twenty-four, and 10.99\% at twenty-four to thirty-six. The second is rate regime. Splitting the pooled window at the 2018/2019 turn separates a rising-rate leg (the 30-year rate moved 4.03\% to 4.64\% over 2018, suppressing refinancing incentive) from a falling-rate leg (4.46\% to 3.72\% over 2019, a refinancing wave). The legs read very differently: the 2018 leg returns 4.70--5.33\% across the out-of-the-money gap grid at the age-$\geq$-12 cut, while the 2019 leg returns 6.91\% and carries the entire 10.99\% age-$\geq$-24 cell. The pooled 6.1\% is therefore a refinancing-inflated composite, not a baseline turnover read, and the 2019 leg is disqualified on both diagnostics at once. The rate-regime half of that disqualification, as stated, compares the two legs at different discount depths; the matched-depth reconciliation below repairs the comparison and reports what survives it, which is less than the sentence just written implies.

What remains is the 2018 rising-rate leg, and its status must be stated precisely, because the paper's off-window anchor now rests on it. Its mature-age seasoning test is \emph{not computable}: the leg contains no out-of-the-money cohort-months at all past the twenty-four-month age cut. A loan that is out of the money while rates are rising is necessarily a recent, low-coupon origination, and therefore young. The 2018 leg is thus clean by rate-regime construction, not by test. Its one computable age contrast (3.78\% to 5.33\% across the twelve-month cut) is confounded by the PSA ramp, on which young loans genuinely turn over less, so it is not evidence about the contrast the leg is missing. The out-of-agency read is that in every cell where a mature contrast exists at all, mature turnover is the higher one, and Section~\ref{sec:identification} reports by how much. A second limitation concerns depth. The gap-$\leq$-two-point cell at which the in-window anchor is measured does exist off-window: the 2018 leg's gap runs to $-2.87$ points, and it holds 35 cohort-months reading 3.440\% annual CPR. Grids that stop at half a point do so because their thresholds were fixed at zero, a quarter point, and half a point in the script header, not because the rate history forbade anything deeper. The true limitation is exposure collapse, and a matched-depth run (\texttt{matched\_depth\_reconciliation}) measures it. Down the 2018 leg's depth ladder the share of its gap-$\leq$-0 exposure that survives runs 100\%, 83.7\%, 72.9\%, 35.9\%, 11.3\%, 3.9\%, and 0.02\%, so the two-point cell carries \$0.282 billion of \$1{,}264.5 billion. That cell is present and worthless. An ex-ante support rule excludes it before any reading is taken: a depth is admissible for a leg only if the leg retains at least 10\% of its gap-$\leq$-0 exposure there. These are the reasons the off-window floor resolves to a \emph{range}, and I report it as one.


\paragraph{Calibration entries for the headline marginal row.} Fannie Mae read of the same selection rule on the pooled off-window period 5.52\%, above the clean band and bracketing the marginal below $+4.3$, its like-for-like Freddie counterpart 5.185\% rather than the 2018 leg's 4.991\%, so 0.337 points of that distance is agency and the rest is period; $+3.9$ to $+7.5$ joint with the band; additive form $+11.2$ (form-conditional hull $+3.5$ to $+13.1$); concave transform $+5.1$; the two \emph{jointly} $+9.2$, an interaction of $-1.47$ points, so these two entries do not add; Ginnie overlay $+4.4$ (cuts the same way as the age-standardized read; now run jointly with it in run \texttt{b5\_joint\_cell}: composed $+2.9$, within $0.1$ of proportional); floor-stability check $+\$45.1$B; baseline-level sweep 75--150 PSA $+0.9$ to $+15.6$ (a convention range, no coverage property; run \texttt{psa\_level\_sweep}); covariate endpoints move this marginal by only $+0.1$ while spanning 8.2 points of level; survival-selection attenuation $+5.1$/$+4.6$/$+3.7$ at $\eta = 0.86/0.73/0.54$, the frailty-identity mapping $\eta = S_{\mathrm{pre}}^{\xi}$ at the measured $S_{\mathrm{pre}} = 40{,}234/75{,}000$ and $\xi = 0.25/0.5/1$ (run \texttt{attenuation\_sensitivity}; a transport sensitivity on the imported coefficient, listed in Table~\ref{tab:assembly} as a bracket in $\xi$ rather than as a correction).

\paragraph{The inference ladder on the floor read.} Section~\ref{sec:identification} quotes the restricted wild-cluster inversion rung of Table~\ref{tab:ladder} as the binding layer; the full ladder, with its degrees of freedom, follows.

\begin{table}[!t]
\centering
\caption{The inference ladder on the binding layer---binding among the layers with a coverage property; the unestimated baseline-level convention spans wider. Every row is one rung on that layer --- the floor read's own sampling error, propagated through the frozen floor-to-marginal grid at the central elasticity $\delta = 6.5\%$ --- so no row's width comes from the elasticity, which is swept separately in Table~\ref{tab:lowband}. All reads are on the 2018 leg's mid-grid anchor, which 31 stratum clusters carry; the last two rungs re-cluster that same read on the six calendar months it populates and on both margins at once. The restricted wild-cluster inversion is the primary construction and the interval this paper quotes, the Webb-weight wild-cluster bootstrap-$t$ being retained as a reported rung; the percentile read is reported demoted rather than as a competing interval. Runs \texttt{floor\_uncertainty} (percentile), \texttt{floor\_inference\_correction} (Rademacher wild-$t$; CR1--CR3 at $t(30)$), \texttt{floor\_inference\_correction\_v2} (Webb primary, Bell--McCaffrey degrees of freedom, restricted inversion), and \texttt{floor\_inference\_correction\_v3} (month and two-way clustering).}
\label{tab:ladder}
\footnotesize
\begin{threeparttable}
\begin{tabular}{@{}llll@{}}
\toprule
Rung & Marginal (pp) & Degrees of freedom & Status \\
\midrule
Percentile cluster bootstrap & $+3.0$ to $+8.0$ & --- (1{,}000-replicate) & demoted; under-covers at 31 clusters \\
CR1 $t$ & $+3.2$ to $+8.3$ & $t(30)$, $G-1$ & no leverage adjustment \\
CR2 $t$ & $+3.0$ to $+8.6$ & $t(30)$, $G-1$ & conventional df \\
CR3 $t$ & $+2.8$ to $+8.8$ & $t(30)$, $G-1$ & conventional df \\
Wild-$t$, Rademacher & $+2.8$ to $+8.7$ & bootstrap, 9{,}999 replications & the committed construction, replayed \\
Wild-$t$, Webb & $+2.9$ to $+8.7$ & bootstrap, 9{,}999 replications & covers 90.8/91.4; below the pre-committed bar \\
CR1 $t$, Bell--McCaffrey & $+2.8$ to $+8.8$ & $t(6.2)$, data-driven & small-sample df \\
CR2 $t$, Bell--McCaffrey & $+2.4$ to $+9.1$ & $t(5.1)$, data-driven & small-sample df \\
CR3 $t$, Bell--McCaffrey & $+2.3$ to $+9.6$ & data-driven, CR3's own & small-sample df \\
Restricted wild-cluster inversion & $+2.3$ to $+9.1$ & 401-point grid inversion & primary; the binding layer (covers 94.3/95.1) \\
CR1 $t$, month clusters & $+3.1$ to $+8.4$ & $t(5)$, $G-1$ & wider on df alone \\
Two-way, stratum $\times$ month & $+2.3$ to $+9.3$ & $t(5)$, $\min(G)-1$ & lower edge censored \\
\bottomrule
\end{tabular}
\end{threeparttable}
\end{table}

\noindent{\footnotesize \emph{Notes to Table~\ref{tab:ladder}.} The Bell--McCaffrey degrees of freedom are data-driven and estimator-specific, so one value cannot serve all three cluster-robust variants: the 6.2 and 5.1 shown are CR1's and CR2's, and CR3's own is smaller still. CR1 applies no leverage adjustment --- it is the narrowest of the three sandwiches --- and is the standard error the wild rows studentize with, so its two rows are the ladder's baseline rather than competing reads. Its Bell--McCaffrey row printing CR3's conventional-df interval is an accident of this leverage profile, not an identity. Unrounded, CR1 at Bell--McCaffrey degrees of freedom runs $[+2.753, +8.796]$ and CR3 at conventional degrees of freedom $[+2.773, +8.774]$ --- two distinct intervals that agree only once each endpoint is taken to one decimal. The effective cluster count is 5.9 by the same Herfindahl-inverse convention as the loan/stratum scheme's 25.8, against 31 nominal clusters. The run's recorded maximum single-cluster leverage is 0.33, in the 2017-vintage, 4.0\% coupon, FICO $740+$, LTV ${\leq}80$ stratum. That is ten times an equal 31-cluster share and 1.9 times the next largest (the same cohort at LTV ${>}80$, 0.17). That concentration is why the wild-cluster correction is applied and the percentile read demoted. It does not rank the wild construction against the Bell--McCaffrey degrees of freedom, and I do not claim it does. The wild rows re-draw cluster-level weights at that same leverage profile; the Bell--McCaffrey rows keep the sandwich and correct the reference distribution's degrees of freedom for it. At 5.9 effective clusters neither correction dominates. So the primary read is now adjudicated, not asserted. Run \texttt{fewcluster\_coverage}'s pre-committed rule --- narrowest construction with Gaussian coverage $\geq 93.0\%$ and heavy-tail coverage $\geq 91.0\%$ --- selects the restricted inversion. That is the artifact's own \texttt{binding\_construction} field, whose $[+2.3, +9.1]$ this paper quotes as the binding layer. The Webb wild-$t$, at 90.8\%/91.4\%, falls below that bar, and its $[+2.9, +8.7]$ is retained as a reported rung, not the quoted interval. The run's own \texttt{landing\_branch} string records L1 (Webb retains), which contradicts the rule its spec froze. The string is sha-pinned and left unedited, and this note is the adjudication. CR2 at Bell--McCaffrey degrees of freedom ($+2.4$ to $+9.1$ points), the other qualifying rung, stands beside the restricted inversion, over-covering and wider. The three rungs printed wider still understate their own width. CR3 at Bell--McCaffrey degrees of freedom, the restricted inversion and the two-way rung have \emph{lower} endpoints at floors above the committed grid's 6.0\% end, and are censored to it. The marginal falls as the floor rises, so all three run below the $+2.3$ printed. Their upper endpoints are interior. The last two rungs re-cluster the same read rather than re-estimating it: the month rung groups the 137 cohort-months by the six calendar months the 2018 leg populates, July to December, and the two-way rung adds the CR1 variance clustered on the 31 strata to the one clustered on those six months and subtracts the one clustered on their 137 intersection cells, the Cameron--Gelbach--Miller construction, carried at $\min(G)-1 = 5$ degrees of freedom. Both come in wider than the CR1 rung they re-cluster, so the month-level shock is priced here rather than dismissed. They widen for different reasons, and only the two-way rung widens because its variance is larger. Month clustering widens the ladder, as I pre-committed it would and for the reason the specification gave, the drop from 31 clusters to six: CR1 goes from 5.20 to 5.33 points of width, both unrounded. Its variance runs the other way, which that specification allowed. The month-clustered CR1 standard error comes in 18\% \emph{smaller} than the stratum-clustered one, so at this cluster unit the common shock adds no variance, and the widening is the six-cluster reference distribution alone. Clustering on both margins at once does add variance---the Cameron--Gelbach--Miller standard error runs 16\% \emph{above} the stratum-clustered one---and that rung is the wider of the two new rows; whether it would be computable at all I deliberately did not predict. My second pre-commitment missed. I expected the month rung's lower marginal edge to be censored at the grid's 6.0\% end, as the two widest stratum rungs already are, and it is not---both endpoints of the printed rung are interior, and of the two new rows only the two-way one censors. Among the month rung's own variants only the unprinted Rademacher one reaches past that edge. Six clusters is small in more than one direction. The month rung's own Bell--McCaffrey degrees of freedom is 3.9---below the four my pre-committed too-few-clusters rule was written to protect, though that rule's own $G < 5$ trigger did not fire. Six clusters admit only $2^{6} = 64$ distinct Rademacher sign vectors, so a Rademacher wild-$t$ tail probability at this cluster unit cannot resolve below $1/64$, which is why the Webb six-point weights are primary here, as they are above.\par}

\section{Behavioral Extension Mechanics}\label{sec:robustness-extensions}

This appendix gives the mechanism behind each of the three behavioral extensions summarized in Section~\ref{sec:abm-results}. As Appendix~\ref{sec:method-frictions} describes, loss aversion steepens the penalty every household faces. On its own that would push simulated mobility below the historically observed 4--5\% baseline turnover floor. To preserve the floor I re-calibrate the mobility-desire parameter upward, nearly tripling it. That recalibration mechanically inflates the model's overall responsiveness to rate gaps. The jump to 54.9\% is therefore an artifact of holding the turnover floor fixed, not a validation of loss aversion as a mechanism. Because that floor is anchored at maximum lock-in, the sign of the recalibration's effect on rate-sensitivity is not fixed. A specification change that eases turnover at the anchor lowers the recalibrated mobility scale, and deflates rate-sensitivity with it: in the real-covariate cross-design of Section~\ref{sec:robustness-crossdesign} the scale falls to 36,086 from the synthetic 43,883. The headline recovery rests on a floor-anchored calibration too, so it inherits that indeterminacy. That is why the ABM's negative verdict is read only alongside the cross-design hedge, never on its own. The second extension, income-scaled curtailment, pulls the estimate down to 45.1\%. Curtailment is a strictly additive drag on prepayment, so I predicted this decline ex ante from a diagnostic correlation between the ABM's over-prediction and disposable-income growth (Appendix~\ref{sec:method-frictions}). The third extension, multi-vintage coupon composition, pulls the estimate down further, to 33.7\%. A lower weighted-average coupon widens the gap to prevailing market rates. My hypothesis was that this should increase predicted trapped liquidity, not reduce it. The inversion suggests the ABM's CPR mechanism is not translating a wider coupon gap into slower prepayment the way the household-choice framework predicts it should. I have not diagnosed the inversion, and its implication cuts in both directions. An aggregate response to the rate gap running backward is equally consistent with ``household choice explains little'' and with ``this specific ABM's gap response is mis-specified.'' So the ABM's recovery share is never quoted in this paper without the cross-design qualifier of Section~\ref{sec:robustness-crossdesign}, and the paradigm attribution inherits that section's hedge. It is also why the paper asks the reader to carry the two-estimator contrast, not the ABM's point estimate. I cannot rule out that a differently specified household model would perform better, but that possibility does not bear on the results reported in Section~\ref{sec:abm}.

The wait-and-see freeze now carries a direct sensitivity run, scored on the frozen Danish-anchor production accounting (trapped \$84.5 billion at the production 150 bp trigger and 20\% share; parity gate $+\$0.02$ billion). Removing the freeze entirely drops trapped liquidity to \$44.6 billion, an isolated contribution of $+\$39.9$ billion, or $+5.2$ points of benchmark. Moving the six-month trigger to 100 or 200 bp spans \$105.9 to \$70.5 billion. The trigger's placement therefore matters at the few-points scale, and its author-calibrated status remains a disclosed limitation. Sweeping the freeze share with the floor anchor retained gives \$65.1 billion at a 10\% share, \$81.5 billion at the externally implied 18.3\%, and \$103.3 billion at 30\%. Re-anchoring the share to the Liebersohn--Rothstein-implied value moves recovery by only $-0.4$ points of benchmark. The calibrated mobility scale and floor CPR are bit-identical at every share, confirming that the freeze cannot compromise the baseline turnover floor. The front-end DTI gate gets the parallel treatment. Shifting the 43\% wall to 36\% or 50\% on the same frozen surface moves ABM trapped liquidity to \$153.3 and \$119.7 billion, from \$84.5 billion at the production 43\%. The turnover floor is re-anchored into the empirical 4--5\% band at every threshold. The gate is therefore material to the ABM's own estimate while never breaching the floor. That sweep is non-monotone in the wall, and the production 43\% sits at its minimum. This invites the reading that the paper's own parameter choice is the one minimizing its own estimate. That reading does not survive decomposition (run \texttt{dti\_nonmonotonicity\_decomposition}). The floor-retention rule re-calibrates the mobility-desire scale only where the fixed-scale floor would leave the 4--5\% band. Across these three thresholds that happens exactly once: at 50\%, where loosening the wall lifts the fixed-scale floor to 5.09\% and the search cuts the scale from 43{,}883 to 41{,}934. Holding the scale at its production value instead, the sweep is monotone (\$153.4, \$84.6 and \$46.4 billion at 36, 43 and 50\%). The 50\% leg's entire rebound is the re-calibration, worth $-\$73.4$ billion. The frozen and re-calibrated legs are bit-identical at 36 and 43\%, where the rule is a no-op, and that is what identifies the channel. The wall is therefore monotone in the direction the mechanism implies, and the published non-monotonicity is the free parameter absorbing the floor constraint. Nothing here moves a headline: the ABM is a stress-test device and the hazard marginal carries no DTI gate.

\section{Duration Quantification: Approximate Portfolio Weighted-Average Life}\label{sec:robustness-wal}

The mechanism the paper quantifies is duration extension, so I express it in maturity units. Table~\ref{tab:wal} is printed with the cap-design arithmetic in Section~\ref{sec:discussion-capdesign}. It reports the SOMA book's approximate weighted-average life (WAL) under scheduled amortization only, under the empirical CPR path, and under each estimator's implied path, at the QT window's start and end. This appendix documents its construction and reads it against each estimator.

Read against the empirical column, each estimator's CPR error maps into effective-maturity error. The ABM's over-predicted prepayment implies a book roughly 3.4 years too short at window start; Path A's under-prediction, a book 1.5 years too long; Path B lands within 0.3 years. The table's final row gives the no-shock world, with prepayments continuing at the 22.81\% mean 2021 realized CPR and implying a 3.4-year WAL. Against it, the empirical 9.4-year WAL is a 6.0-year extension. That duration gap is what the paper's dollar benchmark prices month by month. The two Danish rows read the counterfactual in the same maturity units. Feeding the Danish legs' mean CPRs through the identical calculator gives a rule-only WAL of 9.1 years at the window's start against the U.S. leg's 9.7---a 0.6-year rule-only shortening, with the Danish-level bracketing case at 10.8. The rule-only leg's extension is additionally rate-invariant on the rate-rise side, by the identified moving-flatness (Section~\ref{sec:abm-danish}). I deliberately price extension, not convexity: option-adjusted duration and convexity require a term-structure and prepayment-option model outside this paper's scope. \citet{hanson2014} shows that the convexity so omitted has aggregate rate-dynamic consequences of its own, which is a reason to flag the omission. The term ``convexity-neutral'' is used strictly in this rate-rise sense.

\section{Fifteen-Year MBS Fold-In and Path B Rate-Gap Units}\label{sec:robustness-pipeline}

Two features of the production pipeline set headline figures cited throughout this paper: how fifteen-year MBS are treated, and the units in which Path B applies its rate-gap coefficient. The replication package's ledger catalogues the figures each supersedes.

The first is the fifteen-year MBS fold-in. The SOMA cohort parser buckets 15-year holdings separately by coupon and never merges 15-year with 30-year buckets. Scheduled amortization on 15-year holdings enters both the empirical CPR back-out and the ABM's simulated roll-off. Restricting to 30-year MBS covers 90.6\% of face value. The 15-year book, approximately 9.1\% of holdings, then contributes to total portfolio size but not to cohort-correct amortization or CPR weighting. Folding it in brings coverage to 99.8\% of SOMA face value at a SOMA weighted-average coupon of 2.49\%. Table~\ref{tab:wal}'s note reconciles these coverage components with the audit parse's 90.7\%/9.1\% term split.

I tested, but did not adopt, a fully native treatment in which 15-year loans also draw voluntary CPR from a 15-year-specific surface. It produced out-of-bounds predicted values. The ABM's payment-delta gate is the change in \eqref{eq:annuity}'s monthly payment upon refinancing at the prevailing rate; it drives loss-aversion and mobility-desire responses elsewhere in this paper. It breaks down for seasoned 15-year borrowers, whose same-term replacement payment is nearly flat relative to their current payment. Native 15-year surfaces predicted 32\%--61\% CPR against an empirical range of roughly 5\%--8\%. That was enough to flip the ABM's aggregate trapped-liquidity estimate to a large negative value, far outside the empirically observed range. Production therefore uses a structural-only fold-in. The 15-year cohorts contribute real weights and real scheduled amortization, while voluntary CPR for those cohorts is drawn, in the frozen headline specification, from the same-coupon 30-year surface. This is a scope limitation for the production headline. It does not bind the framework: a 15-year-specific behavioral gate---a pure rate-lock penalty for short-amortization loans---is implemented, and is the basis of the Berger run's calibration-stage freeze (Table~\ref{tab:danish}, note b). Its adoption for the production headline is deferred to the pre-submission freeze.

The second is the units in which Path B applies its rate-gap coefficient $\beta_1$, calibrated from \citet{liebersohn2024} via the survival-function transform described in Section~\ref{sec:pathb} and expressed as a log-hazard-ratio per 100 bp increase in the rate gap. Applied to the gap in decimal units, the coefficient is numerically inert: at a representative QT-era state (a 3.0\% coupon against a 6.8\% market rate) it yields a multiplier of approximately $\times$1.0026, against the intended suppression of approximately $\times$0.77. Production applies it as $\exp(100\,\beta_1\,\tilde g_{i,t})$, with $\tilde g_{i,t}$ the decimal form of \eqref{eq:pathB}'s rate gap, equivalently $\exp(\beta_1 g_{i,t})$ in \eqref{eq:pathB}'s 100 bp units. The gap is negative when the borrower is locked in, so \eqref{eq:beta1}'s positive $\beta_1$ suppresses prepayment. This application is consistent with Path A's convention of a positive coefficient applied to the gap in comparable units. Production also expresses the Danish counterfactual's rate-gap function in rate-equivalent units, so that the buyback discount properly offsets the locked-in spread. That function is live in the mechanism-substitution variant. The production Danish leg's prepay response, under either Berger-recalibrated anchor of Section~\ref{sec:abm-danish}, does not consume that function's classification, which is computed all the same, as Section~\ref{sec:pathb} discloses. Section~\ref{sec:pathb} reports the results under this convention; the superseded \$747 billion (97.7\%) microsimulation is catalogued in the replication package's ledger.

Both are structural or mechanical, not changes in modeling philosophy, and both are additive. Disabling either one---30-year-only filtering, or the decimal-unit $\beta_1$ convention---exactly reproduces the corresponding superseded figures catalogued in the replication package's ledger. I treat that as evidence that each is isolated to its intended code paths.

\section{Seasonalizing the Turnover Floor: Floor Dispersion, and Why No Timing Claim Follows}\label{sec:robustness-seasonalfloor}

The three checks above vary the floor's level, its functional form, and its provenance, holding it constant within the window in each case. Two runs in this paper let the floor vary within the window instead. The cyclical variant of Section~\ref{sec:pathb} indexes it to the rate path; the pre-committed seasonal check taken up here indexes it to the calendar. Both move the marginal, and this check is the one built to isolate why. It asks two questions at once. The first is whether letting the turnover floor vary by calendar month moves the lock-in marginal. The second is whether the resulting monthly path tracks the empirical benchmark well enough to license a timing claim. The run is \texttt{seasonal\_floor\_timing}; its spec, leg definitions, placebo families, and a frozen timing rule were fixed in the script header before any run executed. The answers are asymmetric, and the second is the more consequential. The marginal does move, but not for the reason the design was built to test, and the timing rule is met without being informative.

The calendar profile is a twelve-cell exposure-weighted annualized turnover CPR measured on deeply out-of-the-money cohort-months. I estimate it on the full cohort-month panel (7{,}756 cohort-months, \$74.2 trillion of exposure, pooled annual CPR 4.187\%), not on the QT window. The window-restricted variant differs in four calendar cells and drives nothing here. I state the basis because an earlier module docstring described the profile as in-window and was wrong. The profile peaks in June, with a peak-to-trough ratio of 2.13. I run four legs, each as a paired central and $\beta_1 = 0$ null: the production constant 4.0\% floor, a constant 4.5\% floor, a \emph{shape-only} leg carrying the calendar profile with its exposure-weighted mean pinned at 4.0\%, and the pre-committed 4.5\%-mean seasonal leg. The shape-only leg is the one that isolates the calendar. The pre-committed leg confounds shape with the level move the committed sweep already prices. Every leg re-runs the production microsimulation with only \texttt{INVOLUNTARY\_CPR\_ANNUAL} varied by calendar month. The constant-4.0\% control reproduces the committed production artifacts bitwise: central \$818.530 billion, null \$748.185 billion, marginal $+\$70.345$ billion, $+9.20$ points.

Giving the floor a calendar profile at a held-constant 4.0\% exposure-weighted mean moves the marginal from $+\$70.345$ billion to $+\$65.448$ billion, $+9.20$ to $+8.56$ points: a shape effect of $-\$4.897$ billion. A permutation family disposes of the natural reading of that number. I assign the \emph{same} twelve floor values to scrambled calendar months, so the multiset of values is preserved exactly and only the month labels change. Across fourteen draws, 89.2\% of the shape effect survives on average (69.9\% to 102.8\%); the permuted marginals run \$65.311 to \$66.924 billion against the true \$65.448 billion. Calendar order is very nearly irrelevant to the movement, so the movement is not a calendar effect.

It is instead a Jensen effect operating through the floor's combination rule. Production combines the floor and the voluntary hazard by a hard maximum (\texttt{FLOOR\_MODE = 'max'}, the semantics \eqref{eq:pathB} and the form test above discuss). At the production floor the floor binds in 36.3\% of evaluated loan-months. Under a maximum, dispersion in the floor is one-sided: where the voluntary hazard already exceeds the floor, the low floor cells are truncated away and contribute nothing, while the high cells bind and lift the effective hazard. Any dispersion in the floor therefore raises book-wide prepayment regardless of which month carries which value---exactly what the permutation family measures. A convexity-matched control separates the small remaining piece. The annual-CPR-to-monthly-hazard transform is itself convex, so the twelve-cell profile is equivalent in level to a flat 4.0090\% floor rather than a flat 4.0000\% one. Running that flat leg gives $+\$70.156$ billion, a $-\$0.189$ billion move. Level leakage is thus 3.9\% of the shape effect, and floor dispersion under the maximum rule is the remaining 96.1\%. I report this as a floor-dispersion and functional-form result. It is not evidence that seasonality moves the marginal.

Level and shape do not separate cleanly, and I do not present an additive split as though they did. Measured against the constant-4.0\% control, the level move contributes $-\$12.704$ billion ($-1.66$ points) and the shape move $-\$4.897$ billion ($-0.64$ points); the two together come to $-\$16.040$ billion ($-2.10$ points), leaving an interaction residual of $+\$1.560$ billion ($+0.20$ points). On the simulated monthly paths that residual is larger relative to what it qualifies---mean absolute interaction residual 0.164 points of monthly CPR against mean absolute total movement 0.533 points, so the interaction runs 30.9\% of the total movement. Any level-plus-shape decomposition of this floor must carry the residual; quoting either component alone misstates both.

What the experiment does buy is a narrow bound. Holding the turnover floor flat within the window is a modeling convenience. This run bounds how much of the constant-floor marginal that convenience was responsible for: at most $\$4.897$ billion of the $+\$70.345$ billion in-sample marginal, 7.0\% of it. That same absolute bound is 11.5\% of the $+\$42.6$ billion headline marginal it is measured against. Of even that, 96.1\% is the maximum rule's response to dispersion. The identified lock-in marginal is not an artifact of a flat floor. That is the whole of what the seasonal design establishes, and it establishes it about levels only.

\paragraph{The timing concession.} The design's second purpose was to test whether a seasonalized floor earns a monthly-timing claim, under a rule frozen before the runs: Theil $U_2$ on first differences below one \emph{and} a positive detrended lag-0 correlation against the empirical benchmark. That rule is met as written by both seasonal legs (shape-only $U_2 = 0.995$, $r = +0.016$; pre-committed leg $U_2 = 0.990$, $r = +0.070$), and I report it as met, because it was frozen ex ante and I will not reopen a rule after seeing which way it fell. I nonetheless claim nothing from the pass, because five diagnostics computed afterwards show the rule cannot discriminate. First, the benchmark carries it. The SOMA back-out has four exact-zero months (June 2022, February 2023, April 2024, September 2025), three of them immediately followed by a spike of up to 14.01\% annualized. The zeros are now traced to their mechanism: each is clip-induced, a month whose reported rolloff falls below the scheduled-amortization floor so that the back-out's non-negativity clip pins it at exactly zero, with the displaced mass posting into the following month's reading (unclipped values $-1.59$/$-0.49$/$-0.84$/$-0.84$ points; run \texttt{h1\_zero\_months\_diagnosis}) --- a month-boundary allocation artifact of the reporting series, not a behavioral observation. The obvious alternative --- rebuilding the series from published monthly principal payments, or from per-CUSIP factor changes --- was probed rather than assumed, and neither route is usable. No monthly SOMA principal-payment series is published: the summary endpoint carries holdings levels only, and both monthly endpoints return HTTP 400. The per-CUSIP route does exist but recovers face levels rather than principal payments --- gross face declines run \$2{,}186 billion against a \$652.8 billion net, since face also falls on reinvestment settlement and roll --- so netted over 192 weekly as-of dates it reproduces the committed month-end differencing, \$652.7517 billion against \$652.7517 billion, and returns the same four low months at $-\$2.0$, \$4.5, \$3.7 and \$3.6 billion against a \$16.8 billion window median. Three of the four sit on a week the New York Fed republished unchanged and June 2022 does not, and republication is ordinary rather than diagnostic --- 59 of those 192 weeks and 36 of the 44 months they span carry one --- so what the rebuild establishes is that the alternative source is unavailable and the zeros survive it, not a second mechanism (run \texttt{benchmark\_monthly\_rebuild}). The June 2022 zero opens the window, is followed by an ordinary month, and so contributes one first difference where the others contribute two. Those pairs touch 7 of 41 first differences in total but carry 0.813 of the empirical differenced signal and 0.796 to 0.811 of every leg's $U_2$ numerator. Repairing them by reallocating each spike evenly across its pair, conserving the pair total, flips both seasonal legs from pass to concede and turns the shape-only leg's detrended correlation from $+0.016$ to $-0.134$. Second, a calendar-rotation placebo moves the true profile to the wrong months while holding its amplitude and dispersion exactly; the rule is cleared by 1 of 11 wrong-month rotations for each seasonal leg, and the five-month rotation \emph{beats} the true profile on $U_2$ (0.9934 against 0.9947). Third, 3 of the 14 scrambled calendar orders above clear the rule, and the best of them attains $U_2 = 0.975$, better than either real seasonal leg. Fourth, and decisively, the $\beta_1 = 0$ null clears the rule on both constant-floor legs, the 4.0\% null doing so with $U_2 = 0.9997$ and a detrended correlation of $+0.115$, stronger than either seasonal central leg's. A criterion satisfied by a specification with the lock-in elasticity removed cannot be evidence for the lock-in mechanism. Fifth, the correlations are not distinguishable from zero on 42 months in any case ($p = 0.921$ and $p = 0.661$), and the leg that isolates the mechanism is the fragile one: the shape-only leg's verdict flips on 5 of 42 single-month deletions while the pre-committed leg's flips on none.

Two further diagnostics point the same way without being needed for the verdict. The seasonal legs concentrate 79.7\% (shape-only) and 89.4\% (pre-committed) of their simulated detrended variance in the twelve-month calendar harmonic. The empirical benchmark carries 0.8\% of its own detrended variance there. The model acquires calendar structure the benchmark can neither confirm nor refute. On the committed cross-correlation metric, meanwhile, the seasonal legs peak at lag $+6$ with a negative correlation. That peak sits on the boundary of a 42-month window, and I mark it non-interpretable.

The right reading of this is a finding about the criterion, not a failure of the experiment. The frozen rule was a reasonable ex-ante test. It turns out to have no power against the alternatives that matter: a lumpy benchmark, wrong-month profiles, scrambled calendars, and an elasticity-free null all satisfy it. I therefore restrict every claim in this paper to levels. No estimator here earns a monthly-timing claim. That extends to the seasonalized floor the levels-only concession the production microsimulation's monthly path already carries, where the raw lag-0 correlation does not survive detrending (Section~\ref{sec:pathb}). This restricts what the design identifies. It retracts no level result, and no aggregate, marginal, or bound reported elsewhere in this paper depends on timing.

\paragraph{Carried caveats.} Four qualifications travel with the seasonal floor wherever it is quoted, and none of them is repaired by the runs above. First, the profile's identification is circular in the same way the in-window floor's is: the deep-out-of-the-money filter that defines it captures 98.2\% of QT exposure, so it is close to the identity map on the panel and the profile is not measured off a lock-in-free population. Second, the clean off-window leg that rescues the floor's \emph{level} cannot rescue its \emph{shape}: the 2017--2019 out-of-the-money cells number 335 cohort-months on \$86.2 billion of exposure and record zero prepaid balance in February and April, so no twelve-cell profile is estimable there at all. Third, the seasonal share is not stable across windows, the calendar terms explaining $R^2 = 0.011$ of the full-panel variation against 0.755 in the QT window, and the profile inherits that instability. Fourth, this floor profile peaks in June while the committed \emph{voluntary} seasonality of the estimated cohort-hazard form peaks in March and October. They are different objects and need not agree, but the tension is real and I state it rather than let the two profiles sit unreconciled in the same paper.

\begin{table}[H]
\centering
\footnotesize
\begin{threeparttable}
\caption{Seasonalizing the turnover floor. Each leg re-runs the production microsimulation as a paired central and $\beta_1 = 0$ null with only the turnover floor's calendar profile varied; the marginal is the central-minus-null difference and is basis-invariant. $U_2$ and $r$ are scored on the central leg against the SOMA back-out over the 42-month window. The constant-4.0\% row reproduces the committed production artifacts bitwise. Pre-committed run \texttt{seasonal\_floor\_timing}; artifacts catalogued in the replication package's ledger.}
\label{tab:seasonalfloor}
\begin{tabular}{@{}lrrrrrr@{}}
\toprule
Floor leg & Null (\$B) & Central (\$B) & Marginal (\$B) & Marginal (pp) & $U_2$ (diffs) & $r$ (detr.) \\
\midrule
Constant 4.0\% (production) & 748.185 & 818.530 & $+70.345$ & $+9.20$ & 1.0029 & $-0.203$ \\
Constant 4.5\% & 738.869 & 796.510 & $+57.642$ & $+7.54$ & 1.0023 & $-0.255$ \\
Shape only, mean pinned 4.0\% & 742.625 & 808.073 & $+65.448$ & $+8.56$ & 0.9947 & $+0.016$ \\
Pre-committed seasonal, mean 4.5\% & 729.298 & 783.602 & $+54.305$ & $+7.10$ & 0.9896 & $+0.070$ \\
\midrule
\multicolumn{7}{@{}l}{\textit{Placebo families, against the shape-only leg}} \\
Permutation, 14 scrambled orders\tnote{a} & --- & --- & $+65.311$--$66.924$ & --- & 0.975 best & --- \\
Rotation, 11 wrong-month shifts\tnote{b} & --- & --- & --- & --- & 0.9934 best & --- \\
Convexity-matched flat 4.0090\%\tnote{c} & --- & --- & $+70.156$ & $+9.17$ & 1.0029 & $-0.204$ \\
\bottomrule
\end{tabular}
\begin{tablenotes}[flushleft]\footnotesize
\item[a] The same twelve floor values assigned to scrambled months, multiset preserved exactly. The family retains 89.2\% of the shape effect on average (69.9\%--102.8\%), so the movement is floor dispersion under the maximum combination rule, not calendar timing. Three of the fourteen clear the frozen timing rule, and the best beats both seasonal legs on $U_2$.
\item[b] The true profile moved to the wrong months, amplitude and dispersion held. One of eleven clears the frozen rule for each seasonal leg, and the five-month shift beats the true profile's $U_2 = 0.9947$.
\item[c] The flat floor whose monthly hazard matches the calendar-frequency-weighted mean of the twelve seasonal hazards, reproducing the transform's convexity with zero dispersion. Its $-\$0.189$ billion move is pure level leakage, 3.9\% of the shape effect; the remaining 96.1\% is floor dispersion.
\end{tablenotes}
\end{threeparttable}
\end{table}

\paragraph{The within-run cyclical floor: dispersion versus co-movement.}
The within-run cyclical floor variant of Section~\ref{sec:pathb} sets $\underline{h}_t = 4\% \times (1 + \kappa z_t)$ and sweeps $\kappa$ symmetrically about zero. The decomposition of the resulting grid follows.

That range is mostly not a measure of cyclicality, and I no longer read it as one. A separate pre-committed run (\texttt{floor\_cyclical\_permutation}) reproduces all seven committed $\kappa$ cells bit-exactly and then decomposes them. Scrambling the floor path's time order within each $\kappa$ cell, which preserves the multiset of monthly floor values and their mean exactly, leaves 85.0\% of the range span standing: 72 draws, twelve per non-zero $\kappa$, span $+7.20$ to $+9.12$ points against the committed $+7.09$ to $+9.36$. Switching only the combination rule removes it. Under the competing-risks (additive) form the identical $\kappa$ grid spans $+11.22$ to $+11.27$ points, a span of $0.045$ points, so changing the rule alone eliminates 98.0\% of the max-form span; the additive form moves less across the whole $\kappa$ grid than it moves across the 3--5\% constant-floor level sweep. The mechanism is the one-sided Jensen effect that Appendix~\ref{sec:robustness-seasonalfloor} diagnoses at length for the seasonal floor. Under a hard maximum, dispersion in $\underline{h}_t$ lifts the effective hazard whenever the high cells bind and is truncated away where they do not, regardless of \emph{when} it occurs, so $\kappa$ is a dispersion dial before it is a cyclicality dial. Decomposing the low endpoint confirms it. At $\kappa = +0.5$, which departs $-\$16.09$ billion from production, level equivalence accounts for $-\$2.14$ billion (13.3\%), order-free dispersion for $-\$11.42$ billion (71.0\%), and genuine time-order for $-\$2.53$ billion (15.7\%): 84.3\% of that endpoint is order-free. The high endpoint is not an endpoint at all---it departs from production by $+\$1.23$ billion, $+0.16$ points, at $\kappa = -0.1$.

A cyclicality signal does survive. It is real but small. The true marginal falls outside the entire twelve-draw permutation range at every one of the six non-zero $\kappa$ (0 of 72 draws bracket it). The order components are monotone in $|\kappa|$ and sign-consistent with the co-movement story the concession in Section~\ref{sec:pathb} states: $+\$5.29$, $+\$4.77$, and $+\$2.69$ billion at $\kappa = -0.5$, $-0.25$, $-0.1$, and $-\$1.42$, $-\$1.75$, and $-\$2.53$ billion at $+0.1$, $+0.25$, $+0.5$. A floor rising with rates does lower the marginal, and one falling with rates does raise it. But scrambling the time order strips most of the grid's asymmetry while leaving its level drop intact. The permuted means differ across the sign of $\kappa$ by only $\$0.22$, $\$1.80$, and $\$6.01$ billion at $|\kappa| = 0.1$, $0.25$, and $0.5$, against a true spread of $\$13.82$ billion at $|\kappa| = 0.5$ ($+\$68.08$ billion at $\kappa = -0.5$ against $+\$54.26$ billion at $+0.5$). A two-sided response that survives losing the time order is the signature of dispersion, not of co-movement. The symmetric part of the grid's movement is large and is dispersion; the asymmetric part is small and is cyclicality. What the grid bounds is therefore the marginal's sensitivity to within-run floor variation of this amplitude. Of that, at most \$5.29 billion in any cell is attributable to the direction of co-movement. That the concession's residual is bounded by the level sweep already reported holds a fortiori. One correction to the specification as earlier described: the window mean is pinned at the production 4\% only where the floor's non-negativity clip does not bind. It binds in the two $|\kappa| = 0.5$ stress cells, whose realized window means are 4.0005\% and 4.0828\% (the artifact flags the clip per cell). Of those two, $\kappa = +0.5$ sets the bracket's low end at $+7.09$ points, while $\kappa = -0.5$ sits mid-grid at $+8.90$. The invariant therefore fails in the cell that matters most for the reported range. The convexity-matched flat floor equivalent to the $\kappa = +0.5$ path is 4.0977\%, not 4\%.

\end{document}
